\documentclass{article}
\usepackage{iclr2027_conference,times}
\usepackage[utf8]{inputenc}
\usepackage[T1]{fontenc}
\usepackage[italian]{babel}
\usepackage{amsmath,amssymb,amsfonts,amsthm,mathtools}
\usepackage{booktabs}
\usepackage{xcolor}
\usepackage{tikz-cd}
\usepackage{hyperref}
\usepackage{url}
\usepackage{microtype}
\usepackage{clrscode3e}

\title{Semianelli termodinamici e survey propagation\\
\large Una mappa delle deformazioni di SP e BSP nel $K$-SAT casuale}

\author{Carlo Nicolini\\
Raffaele Marino\\
\texttt{c.nicolini@ipazia.com}}

\newtheorem{proposizione}{Proposizione}
\newtheorem{lemma}[proposizione]{Lemma}
\newtheorem{corollario}[proposizione]{Corollario}
\theoremstyle{definition}
\newtheorem{definizione}[proposizione]{Definizione}
\theoremstyle{remark}
\newtheorem{osservazione}[proposizione]{Osservazione}

\newcommand{\E}{\mathbb{E}}
\newcommand{\R}{\mathbb{R}}
\newcommand{\Sig}{\Sigma}
\newcommand{\Tact}{T_{\mathrm{act}}}
\newcommand{\dd}{\mathrm{d}}
\newcommand{\KL}{\mathrm{KL}}
\newcommand{\Sh}{\mathrm{Sh}}
\newcommand{\Ry}{\mathrm{Ry}}
\newcommand{\Ts}{\mathrm{Ts}}
\newcommand{\lnq}{\ln_q}
\newcommand{\eq}{e_q}
\newcommand{\us}{\partial^{s}_{a} i}
\newcommand{\uu}{\partial^{u}_{a} i}
\DeclareMathOperator*{\argmin}{arg\,min}
\DeclareMathOperator*{\argmax}{arg\,max}

\begin{document}
\maketitle

\begin{abstract}
Nel $K$-SAT casuale, vicino alla soglia di soddisfacibilità, le soluzioni si dividono in un numero esponenziale di cluster.
La survey propagation (SP) conta questi cluster e guida la decimazione.
Il backtracking survey propagation (BSP) aggiunge mosse che liberano variabili già fissate.
Esistono già molte deformazioni di SP: la versione energetica SP-$y$, la versione entropica con parametro di Parisi $m$, la famiglia SP$(\rho)$ di Maneva, Mossel e Wainwright.
La maggior parte del lavoro non propone un nuovo algoritmo.
Organizziamo le deformazioni note in un unico quadro, quello dei semianelli termodinamici di Marcolli e Thorngren, e deriviamo passo per passo tutte le equazioni.
Mostriamo che la somma termodinamica di Shannon descrive due livelli distinti: l'aggregazione locale dei warning, con un'entropia relativa, e il potenziale 1RSB sui cluster.
Mostriamo che l'ordine di Rényi dei pesi dei cluster coincide con il parametro di Parisi $m$ e dà lo spettro multifrattale dei cluster.
Per l'entropia di Tsallis ricaviamo la forma chiusa della somma deformata, il suo limite a temperatura nulla e il suo effetto sul potenziale dei cluster.
Con $q$ fisso la deformazione degenera nei limiti noti.
Con $q-1$ di ordine $1/N$ essa dà un ensemble canonico generalizzato, cioè SP-$y$ con un $y$ efficace che dipende dall'energia.
Infine ricaviamo un'identità esatta per la variazione della complessità quando si fissa o si libera una variabile, e la confrontiamo con i criteri di BSP e di Battaglia, Kolář e Zecchina.
Dall'identità e dalle equazioni di Bellman soft del Soft Actor-Critic ricaviamo una variante di BSP con un look-ahead soft a due passi, e verifichiamo in Lean~4 le affermazioni algebriche su cui si basa.
Un insieme di diagrammi commutativi riassume quali limiti portano a quali algoritmi.
\end{abstract}

\section{Introduzione}
\label{sec:intro}

Il problema di soddisfacibilità booleana casuale ($K$-SAT) è un banco di prova per la meccanica statistica dei sistemi disordinati e per gli algoritmi di ricerca~\citep{mezard2009information}.
Al crescere della densità di clausole $\alpha$ le soluzioni si dividono in cluster, poi la misura si condensa su pochi cluster, infine le soluzioni scompaiono alla soglia $\alpha_s$~\citep{mezard2002analytic,krzakala2007clusters,montanari2008clusters}.
Il metodo della cavità con rottura di simmetria di replica a un passo (1RSB) descrive questa geometria~\citep{mezard2001bethe,mezard2003cavity}.
La survey propagation (SP) è il limite della 1RSB che conta i cluster senza pesarli~\citep{mezard2002analytic,mezard2002random,braunstein2005survey}.
La decimazione guidata da SP fissa le variabili più polarizzate e ricalcola i messaggi sulla formula ridotta.
Il BSP di Marino, Parisi e Ricci-Tersenghi~\citep{marino2016bsp} alterna la decimazione con mosse di backtracking e raggiunge densità più vicine alla soglia $\alpha_s$.

Le deformazioni di SP sono note da tempo.
Mézard e Zecchina~\citep{mezard2002random} introducono un parametro di reweighting $y$ che pesa gli stati con la loro energia, e Battaglia, Kolář e Zecchina~\citep{battaglia2004minimizing} usano la versione SP-$y$ con backtracking per minimizzare l'energia nella fase UNSAT.
Mézard, Palassini e Rivoire~\citep{mezard2005landscape} e Krzakala e coautori~\citep{krzakala2007clusters} pesano i cluster di soluzioni con la loro dimensione elevata al parametro di Parisi $m$.
Braunstein e Zecchina~\citep{braunstein2004surveybp} e Maneva, Mossel e Wainwright~\citep{maneva2007survey} mostrano che SP è belief propagation (BP) su uno spazio esteso, e la famiglia SP$(\rho)$ interpola fra SP e BP.
Wemmenhove e Kappen~\citep{wemmenhove2006finite} studiano SP a temperatura finita.

Marcolli e Thorngren~\citep{marcolli2014thermodynamic} definiscono i semianelli termodinamici.
In essi la somma di due numeri è un problema variazionale con un termine di entropia e una temperatura $T$.
Con l'entropia di Shannon si ottiene l'operazione $-T\log(e^{-x/T}+e^{-y/T})$, che per $T\to 0$ diventa il minimo tropicale.
Con altre entropie (Rényi, Tsallis, entropia relativa) l'operazione perde alcune proprietà algebriche o richiede una deformazione del prodotto.
La domanda di questo lavoro è la seguente.
Quali deformazioni di SP sono esattamente deformazioni di semianello, a quale livello della teoria agiscono, e quali limiti le collegano a SP, a BP, a SP-$y$ e a BSP?

\paragraph{Cosa è noto e cosa è nostro.}
Le equazioni SP, SP-$y$, SP$(\rho)$, il potenziale 1RSB, il criterio di bias di BSP e l'indice di backtracking di Battaglia, Kolář e Zecchina sono risultati noti, e li citiamo come tali.
Anche i semianelli termodinamici e le loro proprietà di associatività sono risultati noti~\citep{marcolli2014thermodynamic}.
I nostri contributi sono osservazioni con dimostrazione, che collegano questi risultati.
\begin{enumerate}
\item Una derivazione passo per passo di SP-$y$ per il $K$-SAT, scritta come somma termodinamica con entropia relativa rispetto alla distribuzione dei warning (Sezione~\ref{sec:spy}).
\item La distinzione fra due livelli: la somma di Shannon sui cluster definisce il potenziale $\Phi(y)$, e da essa segue il reweighting locale $e^{-y\Delta E}$. La temperatura locale e il parametro $y$ quindi non sono indipendenti (Sezione~\ref{sec:spy}).
\item L'identità fra l'ordine di Rényi dei pesi dei cluster e il parametro di Parisi $m$, con la lettura multifrattale e il comportamento vicino ad $\alpha_s$ (Sezione~\ref{sec:qaxis}).
\item La forma chiusa della somma di Tsallis deformata, il suo limite $T\to0$ e l'effetto di una deformazione di Tsallis sul potenziale dei cluster (Sezioni~\ref{sec:semirings} e~\ref{sec:qaxis}).
\item Un'identità esatta per la variazione del funzionale di Bethe della complessità quando si fissa una variabile, con il resto di secondo ordine (Sezione~\ref{sec:response}).
\item Una variante di BSP con un look-ahead soft a due passi, che collega le equazioni di Bellman soft del Soft Actor-Critic~\citep{haarnoja2018soft} all'identità locale, con le affermazioni algebriche verificate in Lean~4 (Sezione~\ref{sec:softq}).
\item Diagrammi commutativi che riassumono i limiti (Sezione~\ref{sec:diagrams}).
\end{enumerate}
Le prove numeriche delle deformazioni fatte finora con il nostro codice sono troppo piccole per distinguere un effetto dal rumore statistico.

\paragraph{Guida alla lettura.}
La Sezione~\ref{sec:notation} definisce il problema, i cluster e i due limiti del potenziale 1RSB.
La Sezione~\ref{sec:spbsp} richiama SP, il suo funzionale di complessità e gli algoritmi SID, BSP e SP-$y$ con backtracking.
La Sezione~\ref{sec:semirings} richiama i semianelli termodinamici e ne deriva le formule che usiamo.
Le Sezioni~\ref{sec:spy}--\ref{sec:response} contengono le derivazioni.
La Sezione~\ref{sec:softq} propone la variante con look-ahead soft.
La Sezione~\ref{sec:diagrams} raccoglie i diagrammi, e la Sezione~\ref{sec:discussion} elenca i problemi aperti.

\section{Problema e notazione}
\label{sec:notation}

\subsection{Il $K$-SAT casuale e il factor graph}

Una formula ha $N$ variabili $x_i\in\{-1,+1\}$, con $i=1,\dots,N$, e $M=\alpha N$ clausole.
Il valore $+1$ indica ``vero''.
Ogni clausola $a$ contiene un insieme $\partial a$ di $K$ variabili distinte.
Per ogni variabile $i\in\partial a$ il numero $s_{a,i}\in\{-1,+1\}$ indica il valore che soddisfa il letterale: il letterale è $x_i$ se $s_{a,i}=+1$ ed è $\neg x_i$ se $s_{a,i}=-1$.
La clausola $a$ è violata se e solo se $x_i=-s_{a,i}$ per ogni $i\in\partial a$.
Nell'ensemble casuale ogni clausola sceglie $K$ variabili a caso in modo uniforme e segni $s_{a,i}$ indipendenti ed equiprobabili.

L'energia è il numero di clausole violate,
\begin{equation}
E(x)=\sum_{a=1}^{M}\ \prod_{i\in\partial a}\frac{1-s_{a,i}x_i}{2}.
\label{eq:energy}
\end{equation}
Le soluzioni sono le configurazioni con $E(x)=0$.
Il factor graph è il grafo bipartito con un nodo per ogni variabile, un nodo per ogni clausola e un lato $(i,a)$ quando $i\in\partial a$.
Indichiamo con $\partial i$ l'insieme delle clausole che contengono $i$, e con
\begin{equation}
\partial i^{\pm}=\{a\in\partial i:\ s_{a,i}=\pm1\}
\end{equation}
le clausole soddisfatte da $x_i=\pm1$.
Per un lato $(i,a)$ servono due insiemi di clausole vicine di $i$ diverse da $a$:
\begin{equation}
\us=\{b\in\partial i\setminus a:\ s_{b,i}=s_{a,i}\},\qquad
\uu=\{b\in\partial i\setminus a:\ s_{b,i}=-s_{a,i}\}.
\label{eq:su-sets}
\end{equation}
Le clausole in $\us$ sono soddisfatte dallo stesso valore che soddisfa $a$.
Le clausole in $\uu$ sono soddisfatte dal valore opposto.
Nel limite $N\to\infty$ a $\alpha$ fisso il factor graph è localmente un albero: un intorno di raggio finito di un nodo non contiene cicli con probabilità che tende a uno~\citep{mezard2009information}.

\subsection{Soglie e cluster}

Per $K=3$ la teoria 1RSB prevede le seguenti soglie~\citep{krzakala2007clusters,mertens2006threshold}.
Oltre $\alpha_d\simeq3{,}86$ le soluzioni si dividono in un numero esponenziale di cluster.
Per $K=3$ la soglia di condensazione coincide con $\alpha_d$, mentre per $K\ge4$ vale $\alpha_d<\alpha_c$.
Oltre $\alpha_c$ un numero sotto-esponenziale di cluster porta quasi tutta la misura uniforme sulle soluzioni.
Oltre $\alpha_s\simeq4{,}2667$ non esistono più soluzioni.
Per $K=4$ i valori sono $\alpha_d\simeq9{,}38$, $\alpha_c\simeq9{,}547$ e $\alpha_s\simeq9{,}931$.

Un cluster è, in modo informale, un insieme di soluzioni collegate da cammini di soluzioni in cui ogni passo cambia $o(N)$ variabili, separato dagli altri cluster.
La definizione precisa usa gli stati puri della misura di Gibbs~\citep{krzakala2007clusters,mezard2009information}.
Una variabile è congelata in un cluster se assume lo stesso valore in tutte le soluzioni del cluster.
Indichiamo i cluster con $C$ e la loro entropia interna per variabile con $s_C=N^{-1}\log|C|$.
Il numero di cluster con entropia interna vicina a $s$ cresce come $e^{N\Sig_s(s)}$.
La funzione $\Sig_s(s)$ è la complessità entropica.

Nella fase UNSAT e nello studio del paesaggio di energia si usano gli stati, cioè gruppi di configurazioni di energia minima locale separati da barriere~\citep{mezard2002random,mezard2003cavity}.
Indichiamo gli stati con $\alpha$ e la loro energia con $E_\alpha$.
Il numero di stati con energia vicina a $Ne$ cresce come $e^{N\Sig_e(e)}$.
La funzione $\Sig_e(e)$ è la complessità energetica.
Nella fase SAT gli stati di energia nulla sono i cluster di soluzioni.

\subsection{Il potenziale 1RSB e i suoi due limiti}

Sia $Z_\alpha(\beta)=\sum_{x\in\alpha}e^{-\beta E(x)}$ la funzione di partizione ristretta allo stato $\alpha$, a temperatura inversa $\beta$.
Il potenziale 1RSB è
\begin{equation}
\Xi(\beta,m)=\sum_\alpha Z_\alpha(\beta)^{m},
\label{eq:xi}
\end{equation}
dove $m$ è il parametro di Parisi~\citep{parisi1980sequence,mezard2001bethe,monasson1995structural}.
Il valore $m=1$ dà la funzione di partizione ordinaria.
Il valore $m=0$ conta gli stati senza pesarli.
Due limiti di~\eqref{eq:xi} sono importanti.

\paragraph{Limite entropico.}
Si restringe la somma ai cluster di soluzioni, cioè si prende $\beta\to\infty$ con $m$ fisso nella fase SAT.
Allora $Z_C=|C|=e^{Ns_C}$, e il metodo di Laplace dà
\begin{equation}
\varphi(m)=\lim_{N\to\infty}\frac1N\log\sum_C|C|^{m}=\sup_{s:\ \Sig_s(s)\ge0}\bigl[\Sig_s(s)+ms\bigr].
\label{eq:phi-m}
\end{equation}
Il vincolo $\Sig_s(s)\ge0$ esprime il fatto che una famiglia con meno di un cluster in media non contribuisce.
Se il supremo è interno, il punto di massimo $s^*(m)$ soddisfa $\Sig_s'(s^*)=-m$, e valgono
\begin{equation}
s^*(m)=\varphi'(m),\qquad \Sig_s\bigl(s^*(m)\bigr)=\varphi(m)-m\varphi'(m).
\label{eq:legendre-m}
\end{equation}
Questo limite è studiato in~\citep{mezard2005landscape,krzakala2007clusters,montanari2008clusters}.

\paragraph{Limite energetico.}
Si prende $\beta\to\infty$ e $m\to0$ con $y=\beta m$ fisso~\citep{mezard2002random,mezard2003cavity}.
Scriviamo $Z_\alpha(\beta)^{m}=\exp[-\beta m F_\alpha(\beta)]$ con $F_\alpha=E_\alpha-\beta^{-1}S_\alpha$, dove $S_\alpha$ è l'entropia interna dello stato.
Allora $\beta mF_\alpha=yE_\alpha-mS_\alpha$.
Il termine $mS_\alpha$ vale $(y/\beta)\,O(N)$ e, diviso per $N$, tende a zero.
Resta il peso $e^{-yE_\alpha}$, e si definisce
\begin{equation}
-y\,\Phi(y)=\lim_{N\to\infty}\frac1N\log\sum_\alpha e^{-yE_\alpha}=\sup_{e}\bigl[\Sig_e(e)-ye\bigr].
\label{eq:Phi-y}
\end{equation}
La funzione $\Phi(y)$ ha il ruolo di un'energia libera a ``temperatura'' $1/y$.
Poniamo $G(y)=-y\Phi(y)$.
Il teorema dell'inviluppo dà $G'(y)=-e^*(y)$, quindi
\begin{equation}
e^*(y)=\frac{\partial\,(y\Phi)}{\partial y},\qquad
\Sig_e\bigl(e^*(y)\bigr)=G(y)+y\,e^*(y)=y^2\,\frac{\partial\Phi}{\partial y}.
\label{eq:legendre-y}
\end{equation}
Poiché $\Phi(y)=y^{-1}\inf_e[ye-\Sig_e(e)]$, per ogni $y>0$ vale $\Phi(y)\le e_{\mathrm{gs}}$, dove $e_{\mathrm{gs}}$ è l'energia per variabile degli stati fondamentali, che hanno $\Sig_e(e_{\mathrm{gs}})=0$.
L'uguaglianza vale nel punto $y^*$ in cui $\partial\Phi/\partial y=0$, cioè in cui $\Sig_e(e^*)=0$.
Nella fase UNSAT la stima dell'energia fondamentale è quindi $e_{\mathrm{gs}}=\max_y\Phi(y)$~\citep{mezard2002random}.

\paragraph{Il vertice comune.}
Nella fase SAT si ha $e_{\mathrm{gs}}=0$.
Per $y\to\infty$ la~\eqref{eq:Phi-y} dà $G(y)\to\Sig_e(0)$, il logaritmo del numero di stati di energia nulla, cioè di cluster di soluzioni.
Per $m\to0$ la~\eqref{eq:phi-m} dà $\varphi(0)=\sup_s\Sig_s(s)$, cioè il logaritmo del numero totale di cluster.
I due limiti portano allo stesso oggetto per strade diverse.
Al livello delle equazioni di cavità la corrispondenza è la seguente: il limite $m\to0$ delle equazioni entropiche, ristretto ai messaggi di BP con valori in $\{0,1\}$, dà le equazioni SP~\citep{krzakala2007clusters,montanari2008clusters,mezard2009information}.
I cluster senza variabili congelate non producono warning e sono invisibili a SP.
Per questo motivo la soglia in cui SP ha una soluzione non banale non coincide in generale con $\alpha_d$~\citep{krzakala2007clusters}.

\subsection{Dequantizzazione di Maslov}

Le formule~\eqref{eq:phi-m} e~\eqref{eq:Phi-y} sono entrambe esempi dello stesso limite.
Per $h>0$ si definisce $a\oplus_h b=h\log(e^{a/h}+e^{b/h})$.
Per $h\to0$ si ha $a\oplus_h b\to\max(a,b)$, e il semianello $(\R\cup\{-\infty\},\oplus_h,+)$ tende al semianello tropicale $(\max,+)$.
Questo limite si chiama dequantizzazione di Maslov~\citep{litvinov2007maslov}.
Una somma di $e^{Ng(e)}$ su una famiglia di stati è una somma $\oplus_h$ con $h=1/N$, a meno del fattore di scala.
Il metodo di Laplace, che porta alle trasformate di Legendre~\eqref{eq:phi-m} e~\eqref{eq:Phi-y}, è quindi una dequantizzazione di Maslov con parametro $1/N$.
Nel seguito incontreremo un secondo parametro di dequantizzazione, la temperatura $T=1/y$, e i due limiti non vanno confusi.

\section{Survey propagation, complessità e BSP}
\label{sec:spbsp}

Questa sezione richiama risultati noti.
Fissiamo la notazione che useremo nelle derivazioni.

\subsection{Warning e survey}

Consideriamo un cluster e il lato $(i,a)$.
Il warning $u_{a\to i}\in\{0,1\}$ vale $1$ se, nel cluster, tutte le variabili $j\in\partial a\setminus i$ sono congelate al valore $-s_{a,j}$.
In questo caso la clausola $a$ obbliga $x_i=s_{a,i}$.
La survey $\eta_{a\to i}\in[0,1]$ è la frazione di cluster in cui $u_{a\to i}=1$.
Il messaggio $\nu_{j\to a}\in[0,1]$ è la frazione di cluster in cui $j$ è congelata dalle clausole diverse da $a$ al valore $-s_{a,j}$, cioè al valore che non soddisfa $a$~\citep{mezard2002analytic,braunstein2005survey}.

Assumiamo che, in assenza di $a$, le variabili $j\in\partial a$ siano indipendenti.
Questa ipotesi è esatta su un albero ed è l'ipotesi di cavità sul grafo casuale.
Il warning $u_{a\to i}$ vale $1$ se e solo se ogni $j\in\partial a\setminus i$ è congelata contro $a$, quindi
\begin{equation}
\eta_{a\to i}=\prod_{j\in\partial a\setminus i}\nu_{j\to a}.
\label{eq:sp-clause}
\end{equation}
Per la variabile $j$ e la clausola $a$ usiamo gli insiemi~\eqref{eq:su-sets}.
La variabile $j$ riceve warning da ogni $b\in\partial j\setminus a$ in modo indipendente, con probabilità $\eta_{b\to j}$.
Distinguiamo tre casi, secondo la provenienza dei warning.
\begin{align}
\Pi^u_{j\to a}&=\Bigl[1-\prod_{b\in\partial^u_a j}(1-\eta_{b\to j})\Bigr]\prod_{b\in\partial^s_a j}(1-\eta_{b\to j}),\label{eq:Pi-u}\\
\Pi^s_{j\to a}&=\Bigl[1-\prod_{b\in\partial^s_a j}(1-\eta_{b\to j})\Bigr]\prod_{b\in\partial^u_a j}(1-\eta_{b\to j}),\label{eq:Pi-s}\\
\Pi^0_{j\to a}&=\prod_{b\in\partial j\setminus a}(1-\eta_{b\to j}).\label{eq:Pi-0}
\end{align}
Il caso $u$ indica almeno un warning verso $-s_{a,j}$ e nessuno verso $s_{a,j}$.
Il caso $s$ indica il contrario.
Il caso $0$ indica nessun warning.
Il caso con warning da entrambi i lati è una contraddizione, e nel conteggio dei cluster di soluzioni ha peso zero.
Normalizzando sui casi senza contraddizione si ottiene
\begin{equation}
\nu_{j\to a}=\frac{\Pi^u_{j\to a}}{\Pi^u_{j\to a}+\Pi^s_{j\to a}+\Pi^0_{j\to a}}.
\label{eq:sp-var}
\end{equation}
Le~\eqref{eq:sp-clause} e~\eqref{eq:sp-var} sono le equazioni SP.

\paragraph{Grandezze di sito.}
Per ogni variabile $i$ poniamo
\begin{equation}
\pi_i^{\pm}=1-\prod_{b\in\partial i^{\pm}}(1-\eta_{b\to i}),
\qquad
w_i^{\pm}=\frac{\pi_i^{\pm}(1-\pi_i^{\mp})}{1-\pi_i^+\pi_i^-},
\qquad
w_i^{0}=1-w_i^+-w_i^-.
\label{eq:w}
\end{equation}
Il numero $\pi_i^{+}$ è la probabilità che almeno una clausola spinga $i$ verso $+1$.
Il numero $w_i^{\pm}$ è la frazione di cluster in cui $i$ è congelata a $\pm1$, e $w_i^0$ è la frazione in cui $i$ non è congelata.
Nella notazione di BSP~\citep{marino2016bsp} si ha $\hat m_{a\to i}=\eta_{a\to i}$ e $m_{i\to a}=\nu_{i\to a}$.
La~\eqref{eq:sp-var} coincide con l'equazione~(2) di BSP, perché $\Pi^u+\Pi^s+\Pi^0=1-\pi^s\pi^u$, dove $\pi^s$ e $\pi^u$ sono le probabilità di almeno un warning dai due lati.

\subsection{SP come belief propagation su uno spazio esteso}

Braunstein e Zecchina~\citep{braunstein2004surveybp} e Maneva, Mossel e Wainwright (MMW)~\citep{maneva2007survey} mostrano che SP è BP su un modello con un terzo valore $*$.
Un assegnamento parziale è $x\in\{-1,+1,*\}^N$.
Esso è non valido per la clausola $a$ in due casi: tutte le variabili di $a$ non soddisfano $a$, oppure tutte tranne una non soddisfano $a$ e quella rimasta vale $*$.
Una variabile con valore $\pm1$ è vincolata se è l'unica che soddisfa una clausola, ed è libera altrimenti.
Siano $n_*(x)$ e $n_o(x)$ il numero di variabili con valore $*$ e il numero di variabili libere.
MMW definiscono sugli assegnamenti validi il peso
\begin{equation}
W(x)=\omega_o^{\,n_o(x)}\,\omega_*^{\,n_*(x)}.
\label{eq:mmw}
\end{equation}
Il Teorema~3 di MMW afferma che BP su~\eqref{eq:mmw} con $\omega_*=\rho$ e $\omega_o=1-\rho$ coincide con l'algoritmo SP$(\rho)$.
Il punto $(\omega_o,\omega_*)=(0,1)$ dà SP.
Il punto $(1,0)$ dà BP sulla misura uniforme delle soluzioni.
Per $(0,1)$ i messaggi di BP hanno tre componenti, ma due di esse sono uguali, e la dinamica si chiude su un solo numero per lato, la survey.
Nella Sezione~\ref{sec:qaxis} useremo la condizione precisa per cui questa riduzione vale.

\subsection{Il funzionale di complessità}

Definiamo il funzionale
\begin{equation}
\Sig[\eta]=\sum_a F_a+\sum_i F_i-\sum_{(i,a)}F_{ia},
\label{eq:bethe-sp}
\end{equation}
con i termini
\begin{equation}
F_a=\log\Bigl(1-\prod_{j\in\partial a}\nu_{j\to a}\Bigr),\qquad
F_i=\log\bigl(1-\pi_i^+\pi_i^-\bigr),\qquad
F_{ia}=\log\bigl(1-\nu_{i\to a}\,\eta_{a\to i}\bigr).
\label{eq:bethe-terms}
\end{equation}
Qui le survey $\eta$ sono le variabili indipendenti, e i messaggi $\nu$ sono funzioni di $\eta$ tramite la~\eqref{eq:sp-var}.
Ogni termine è il logaritmo della probabilità che l'aggiunta di un nodo o di un lato non produca contraddizioni.
La Sezione~\ref{sec:spy} ricava questi termini come limite $y\to\infty$ dei termini di SP-$y$.
In un punto fisso vale $\nu_{i\to a}\eta_{a\to i}=\prod_{j\in\partial a}\nu_{j\to a}$ per la~\eqref{eq:sp-clause}, quindi $F_{ia}=F_a$ per ogni $i\in\partial a$.
Il funzionale diventa
\begin{equation}
\Sig=\sum_i\log\bigl(1-\pi_i^+\pi_i^-\bigr)+\sum_a(1-K_a)\log\Bigl(1-\prod_{j\in\partial a}\nu_{j\to a}\Bigr),
\label{eq:sigma-bsp}
\end{equation}
dove $K_a=|\partial a|$.
Questa è la forma usata da BSP~\citep[eq.~(4)--(5)]{marino2016bsp}.

\begin{lemma}[Stazionarietà]
\label{lem:stationary}
Il funzionale~\eqref{eq:bethe-sp} è stazionario rispetto a tutte le survey $\eta$ in ogni punto fisso di SP.
La forma compatta~\eqref{eq:sigma-bsp}, letta come funzione delle $\eta$, in generale non lo è.
\end{lemma}

\begin{proof}
Fissiamo un lato $(k,b)$ e supponiamo $b\in\partial k^+$; il caso $b\in\partial k^-$ è simmetrico.
La survey $\eta_{b\to k}$ compare in tre punti: nel termine di sito $F_k$, nel termine di lato $F_{kb}$, e nei messaggi $\nu_{k\to a}$ con $a\in\partial k\setminus b$.

Primo, i messaggi $\nu_{k\to a}$ compaiono solo in $F_a$ e in $F_{ka}$.
Le derivate sono
\begin{equation*}
\frac{\partial F_a}{\partial\nu_{k\to a}}=-\frac{\prod_{j\in\partial a\setminus k}\nu_{j\to a}}{1-\prod_{j\in\partial a}\nu_{j\to a}},\qquad
\frac{\partial F_{ka}}{\partial\nu_{k\to a}}=-\frac{\eta_{a\to k}}{1-\nu_{k\to a}\eta_{a\to k}}.
\end{equation*}
Se vale la~\eqref{eq:sp-clause}, le due derivate sono uguali, e il contributo tramite $\nu_{k\to a}$ si annulla.

Secondo, poniamo $A=\prod_{c\in\partial k^+\setminus b}(1-\eta_{c\to k})$.
Allora $\pi_k^+=1-A(1-\eta_{b\to k})$, e
\begin{equation*}
\frac{\partial F_k}{\partial\eta_{b\to k}}=-\frac{\pi_k^-A}{1-\pi_k^+\pi_k^-}.
\end{equation*}
Per il lato $(k,b)$ si ha $\partial^s_bk=\partial k^+\setminus b$ e $\partial^u_bk=\partial k^-$.
Le~\eqref{eq:Pi-u}--\eqref{eq:Pi-0} danno $\Pi^u=\pi_k^-A$, $\Pi^s=(1-A)(1-\pi_k^-)$, $\Pi^0=A(1-\pi_k^-)$, con somma $1-\pi_k^-(1-A)$.
Quindi $\nu_{k\to b}=\pi_k^-A/[1-\pi_k^-(1-A)]$.
Un calcolo diretto dà
\begin{equation*}
1-\nu_{k\to b}\eta_{b\to k}=\frac{1-\pi_k^-\bigl(1-A(1-\eta_{b\to k})\bigr)}{1-\pi_k^-(1-A)}=\frac{1-\pi_k^+\pi_k^-}{1-\pi_k^-(1-A)},
\end{equation*}
e quindi
\begin{equation*}
\frac{\partial F_{kb}}{\partial\eta_{b\to k}}=-\frac{\nu_{k\to b}}{1-\nu_{k\to b}\eta_{b\to k}}=-\frac{\pi_k^-A}{1-\pi_k^+\pi_k^-}=\frac{\partial F_k}{\partial\eta_{b\to k}}.
\end{equation*}
Questa identità vale per ogni valore delle survey.
Il termine di sito e il termine di lato entrano in~\eqref{eq:bethe-sp} con segni opposti, quindi il loro contributo si annulla.
La derivata totale è zero.

Nella forma~\eqref{eq:sigma-bsp} i termini di lato $F_{ia}$ sono stati sostituiti da $K_a$ copie del termine di clausola.
Restano due residui rispetto al funzionale di Bethe.
Primo, il termine di sito $F_k$ non è più cancellato da $-F_{kb}$, e $\partial F_k/\partial\eta_{b\to k}$ non è zero in generale.
Secondo, il contributo tramite $\nu_{k\to a}$ ha coefficiente $(1-K_a)\,\partial F_a/\partial\nu_{k\to a}$ invece di zero; per $K_a\neq1$ e $\eta_{a\to k}>0$ questo fattore è non nullo.
Il primo residuo basta già a distruggere la stazionarietà, anche nel caso $K_a=1$ in cui il secondo si annulla.
\end{proof}

Il lemma è un caso particolare di un fatto generale: i punti fissi di BP sono punti stazionari dell'energia libera di Bethe~\citep{yedidia2005constructing}.
Qui lo applichiamo al modello esteso di MMW, e la dimostrazione diretta mostra quali identità servono.
Useremo il lemma nella Sezione~\ref{sec:response}.

\subsection{Decimazione e backtracking}

La decimazione guidata dalle survey (SID)~\citep{mezard2002analytic,braunstein2005survey} alterna due passi.
Prima si risolvono le equazioni SP sulla formula corrente.
Poi si fissa una frazione $f$ delle variabili, quelle con il bias più grande, al loro valore più probabile.
BSP~\citep{marino2016bsp} usa il bias
\begin{equation}
b_i=1-\min(w_i^+,w_i^-).
\label{eq:bias}
\end{equation}
Secondo BSP, se si fissa $x_i$ al valore più probabile, il numero di cluster viene moltiplicato per $b_i$.
La Proposizione~\ref{prop:local} della Sezione~\ref{sec:response} dà la forma esatta di questa affermazione.
Dopo ogni passo si eliminano le clausole soddisfatte, si accorciano le clausole con un letterale falso, e si risolvono di nuovo le equazioni.
Se i messaggi diventano tutti nulli, la formula residua viene passata a WalkSat.

BSP aggiunge il passo di backtracking.
Il bias~\eqref{eq:bias} si calcola anche per una variabile già fissata, con le survey che essa riceve nel punto fisso corrente.
Con una regola stocastica governata dal rapporto $r\in[0,1)$ fra passi di backtracking e passi di decimazione, BSP libera una frazione $f$ delle variabili fissate con il bias più piccolo.
Per $r=0$ si ritrova SID.
Ogni variabile viene assegnata in media $1/(1-r)$ volte, e il costo cresce dello stesso fattore.
BSP usa $f=10^{-3}$.

Battaglia, Kolář e Zecchina~\citep{battaglia2004minimizing} avevano già proposto un backtracking per SP-$y$ (algoritmo SP-Y).
Essi fissano la variabile al valore suggerito da $b_{\mathrm{fix}}(j)=|W_j^+-W_j^-|$, dove $W_j^{\pm}$ sono le probabilità che il campo locale sia positivo o negativo.
Per le variabili fissate usano l'indice con segno
\begin{equation}
b_{\mathrm{back}}(j)=-s_j\bigl(W_j^+-W_j^-\bigr),
\label{eq:bkz-index}
\end{equation}
dove $s_j$ è il valore assegnato.
L'indice è grande quando il campo si è rovesciato contro l'assegnazione.
Anche SP-Y usa un rapporto $r$ fra mosse di rilascio e mosse di fissaggio.
Nel limite $y\to\infty$ si ha $W_j^\pm=w_j^\pm$.

\section{Semianelli termodinamici}
\label{sec:semirings}

Questa sezione richiama le definizioni di Marcolli e Thorngren (MT)~\citep{marcolli2014thermodynamic}.
Le Proposizioni~\ref{prop:shannon} e~\ref{prop:kl} sono note e le dimostriamo per completezza.
La Proposizione~\ref{prop:tsallis} è una derivazione nostra, che usa la distribuzione di Curado e Tsallis~\citep{curado1991generalized}.
Lavoriamo nelle coordinate min-plus, in cui la ``somma'' tropicale è il minimo e il ``prodotto'' è la somma ordinaria.
In queste coordinate un numero $x$ ha il ruolo di un'energia.

\subsection{Definizione}

Sia $\Delta_n=\{p\in[0,1]^n:\ \sum_ip_i=1\}$ il simplesso delle distribuzioni su $n$ elementi.
Sia $S$ una funzione di entropia su $\Delta_n$, e sia $T\ge0$ una temperatura.

\begin{definizione}[MT, Def.~4.1]
\label{def:thermo}
La somma termodinamica di $x_1,\dots,x_n\in\R\cup\{+\infty\}$ è
\begin{equation}
\bigoplus_{S,T}\,(x_1,\dots,x_n)=\min_{p\in\Delta_n}\Bigl[\sum_{i=1}^n p_ix_i-T\,S(p)\Bigr].
\label{eq:thermo-sum}
\end{equation}
Nel caso binario si scrive $x\oplus_{S,T}y=\min_{p\in[0,1]}[px+(1-p)y-TS(p)]$.
\end{definizione}

La struttura $(\R\cup\{+\infty\},\oplus_{S,T},+)$ è il semianello termodinamico.
Per $T=0$ la~\eqref{eq:thermo-sum} è il minimo, e si ottiene il semianello tropicale.
Il termine $\sum_ip_ix_i$ è un'energia media e il termine $-TS(p)$ è il costo entropico, quindi la~\eqref{eq:thermo-sum} è un'energia libera.

\begin{osservazione}[Proprietà algebriche, MT Teor.~3.1 e 4.2]
L'operazione $\oplus_{S,T}$ è commutativa se e solo se $S(p)=S(1-p)$.
L'elemento $+\infty$ è neutro se e solo se $S(0)=S(1)=0$.
L'operazione è associativa se e solo se $S$ soddisfa la relazione di ricorsività
\begin{equation}
S(p_1)+(1-p_1)\,S\!\Bigl(\frac{p_2}{1-p_1}\Bigr)=S(p_1+p_2)+(p_1+p_2)\,S\!\Bigl(\frac{p_1}{p_1+p_2}\Bigr).
\label{eq:assoc}
\end{equation}
A meno di una costante moltiplicativa, l'unica entropia che soddisfa le tre condizioni è quella di Shannon, $\Sh(p)=-p\log p-(1-p)\log(1-p)$.
\end{osservazione}

L'associatività conta per il message passing.
BP e SP aggregano molti messaggi in ingresso con la stessa operazione.
Se l'operazione non è associativa, il risultato dipende dall'ordine in cui si combinano i messaggi.
La legge distributiva generalizzata di Aji e McEliece~\citep{aji2000generalized} richiede un semianello commutativo, cioè una somma associativa e commutativa e un prodotto che distribuisce sulla somma.

\subsection{Shannon: la somma di Gibbs--Maslov}

\begin{proposizione}[MT, Prop.~4.3]
\label{prop:shannon}
Con l'entropia di Shannon $S=H(p)=-\sum_ip_i\log p_i$ e $T>0$ vale
\begin{equation}
\bigoplus_{\Sh,T}(x_1,\dots,x_n)=-T\log\sum_{i=1}^ne^{-x_i/T},
\label{eq:shannon-sum}
\end{equation}
e il minimo è raggiunto nella distribuzione di Gibbs $p_i^*=e^{-x_i/T}/\sum_je^{-x_j/T}$.
\end{proposizione}

\begin{proof}
Poniamo $Z=\sum_je^{-x_j/T}$ e $g_i=e^{-x_i/T}/Z$.
Allora $x_i=-T\log g_i-T\log Z$.
Sostituendo si ottiene, per ogni $p\in\Delta_n$,
\begin{equation*}
\sum_ip_ix_i-TH(p)=-T\log Z+T\sum_ip_i\log\frac{p_i}{g_i}=-T\log Z+T\,\KL(p\,\|\,g).
\end{equation*}
L'entropia relativa $\KL(p\|g)$ è non negativa e vale zero se e solo se $p=g$.
Quindi il minimo vale $-T\log Z$ ed è raggiunto in $p=g$.
\end{proof}

Questa identità è il principio variazionale di Gibbs.
La mappa $x\mapsto e^{-x/T}$ trasforma $\oplus_{\Sh,T}$ nella somma ordinaria e $+$ nel prodotto ordinario.
Quindi, a $T$ fisso, il semianello di Shannon è isomorfo al semianello $(\R_{\ge0},+,\times)$ di BP.
A $T$ fisso esso non contiene nuova algebra: $T$ è una scala.
L'informazione nuova sta nella famiglia in $T$ e nel limite $T\to0$, in cui la~\eqref{eq:shannon-sum} tende a $\min_ix_i$.
Il parametro $T$ coincide con il parametro $h$ della dequantizzazione di Maslov~\citep{litvinov2007maslov}.
Per questa ragione il semianello di Shannon si chiama anche semianello di Gibbs--Maslov.

\subsection{Entropia relativa}

\begin{proposizione}
\label{prop:kl}
Sia $\pi\in\Delta_n$ con $\pi_i>0$.
Con $S(p)=-\KL(p\|\pi)$ e $T>0$ vale
\begin{equation}
\bigoplus_{\KL(\cdot\|\pi),T}(x_1,\dots,x_n)=\min_{p\in\Delta_n}\Bigl[\sum_ip_ix_i+T\,\KL(p\,\|\,\pi)\Bigr]=-T\log\sum_{i=1}^n\pi_i\,e^{-x_i/T}.
\label{eq:kl-sum}
\end{equation}
\end{proposizione}

\begin{proof}
Si ripete la dimostrazione precedente con $g_i=\pi_ie^{-x_i/T}/Z_\pi$ e $Z_\pi=\sum_j\pi_je^{-x_j/T}$.
Si ottiene $\sum_ip_ix_i+T\KL(p\|\pi)=-T\log Z_\pi+T\KL(p\|g)$.
\end{proof}

MT studiano il caso binario nella Sezione~8.
Con $\pi$ uniforme si ha $\KL(p\|\pi)=\log n-H(p)$, e la~\eqref{eq:kl-sum} è la~\eqref{eq:shannon-sum} più la costante $T\log n$.
Con $\pi$ non uniforme l'operazione non è commutativa, perché il peso $\pi_i$ è legato alla posizione $i$.
La~\eqref{eq:kl-sum} è la forma che compare nelle equazioni di cavità: la distribuzione $\pi$ è la probabilità a priori delle configurazioni dei messaggi in ingresso, e $x_i$ è il costo in energia della configurazione $i$ (Sezione~\ref{sec:spy}).

\subsection{Rényi}

L'entropia di Rényi di ordine $\alpha>0$, $\alpha\ne1$, è
\begin{equation}
\Ry_\alpha(p)=\frac{1}{1-\alpha}\log\sum_ip_i^{\alpha},
\label{eq:renyi}
\end{equation}
e tende a $H(p)$ per $\alpha\to1$~\citep{renyi1961measures}.
È additiva su sistemi indipendenti ma non soddisfa la~\eqref{eq:assoc}.
Quindi $\oplus_{\Ry_\alpha,T}$ è commutativa ma non associativa.
Il Lemma~6.1 di MT mostra che il difetto di associatività si corregge con la permutazione $(p_1,p_2,p_3)\mapsto(p_3,p_2,p_1)$.
Per $T\to0$ anche $\oplus_{\Ry_\alpha,T}$ tende al minimo.
Nella Sezione~\ref{sec:qaxis} l'entropia di Rényi non entra come operazione di messaggio ma come misura dei pesi dei cluster.

\subsection{Tsallis e la somma deformata}

L'entropia di Tsallis di indice $q>0$, $q\ne1$, è
\begin{equation}
S_q(p)=\frac{1-\sum_ip_i^{\,q}}{q-1},
\label{eq:tsallis}
\end{equation}
e tende a $H(p)$ per $q\to1$~\citep{tsallis1988possible}.
Non è additiva su sistemi indipendenti: $S_q(A\star B)=S_q(A)+S_q(B)+(1-q)S_q(A)S_q(B)$.
Anche $\oplus_{S_q,T}$ non è associativa.
MT (Sezione~7.1, eq.~(7.4)) introducono un'operazione deformata.
Nelle coordinate min-plus essa si scrive
\begin{equation}
x\oplus_{S,T,q}y=\min_{p\in[0,1]}\bigl[p^{q}x+(1-p)^{q}y-T\,S(p)\bigr].
\label{eq:tsallis-op}
\end{equation}
Il Teorema~7.2 di MT afferma che l'operazione~\eqref{eq:tsallis-op} è associativa se e solo se $S$ soddisfa la versione deformata della~\eqref{eq:assoc}, in cui i fattori $(1-p_1)$ e $(p_1+p_2)$ sono elevati alla $q$.
Il Teorema~7.1 afferma che l'entropia di Tsallis è l'unica entropia commutativa, con la proprietà di identità, che soddisfa questa condizione.
La deformazione sposta la potenza $q$ dal termine di entropia al termine di energia: l'energia media diventa $U_q=\sum_ip_i^{q}x_i$.
Questa è la ``seconda scelta'' di Curado e Tsallis~\citep{curado1991generalized}.

\begin{proposizione}[Forma chiusa della somma di Tsallis]
\label{prop:tsallis}
Siano $x_i\ge0$ e $T>0$.
Poniamo $\lnq u=(u^{1-q}-1)/(1-q)$ ed $\eq(u)=[1+(1-q)u]_+^{1/(1-q)}$, dove $[z]_+=\max(z,0)$.
Allora
\begin{equation}
F_q(x;T)\equiv\min_{p\in\Delta_n}\Bigl[\sum_ip_i^{\,q}x_i-T\,S_q(p)\Bigr]=-T\,\lnq\sum_{i=1}^n\eq(-x_i/T),
\label{eq:tsallis-closed}
\end{equation}
e il minimo è raggiunto in $p_i^*\propto\eq(-x_i/T)$.
Per $q\to1$ la~\eqref{eq:tsallis-closed} tende alla~\eqref{eq:shannon-sum}.
\end{proposizione}

\begin{proof}
Riscriviamo la funzione obiettivo come
\begin{equation*}
\mathcal F(p)=\sum_ip_i^{\,q}x_i+\frac{T}{q-1}\Bigl(\sum_ip_i^{\,q}-1\Bigr)=\frac{T}{q-1}\Bigl[\sum_ip_i^{\,q}\,u_i-1\Bigr],\qquad u_i=1+\frac{(q-1)x_i}{T}.
\end{equation*}
Caso $q>1$.
Si ha $u_i>0$, e $\mathcal F$ è strettamente convessa, perché $p\mapsto p^q$ è strettamente convessa e il coefficiente $u_iT/(q-1)$ è positivo.
Il punto stazionario del lagrangiano $\mathcal F(p)+\lambda(\sum_ip_i-1)$ è quindi il minimo globale.
La condizione di stazionarietà è $q\,p_i^{\,q-1}u_iT/(q-1)=-\lambda$, quindi $p_i\propto u_i^{-1/(q-1)}=\eq(-x_i/T)$.
Poniamo $Z_q=\sum_ju_j^{-1/(q-1)}$, così che $p_i=u_i^{-1/(q-1)}/Z_q$.
Allora $p_i^{\,q}u_i=u_i^{-1/(q-1)}Z_q^{-q}$, e $\sum_ip_i^{\,q}u_i=Z_q^{1-q}$.
Il valore minimo è $T(Z_q^{1-q}-1)/(q-1)=-T\lnq Z_q$.

Caso $0<q<1$.
Si ha $\mathcal F(p)=\frac{T}{1-q}\bigl[1-\sum_ip_i^{\,q}u_i\bigr]$.
Se $u_i\le0$, il termine $-p_i^{\,q}u_i$ è non negativo e il minimo pone $p_i=0$.
Questo è il taglio (cutoff) della distribuzione di Tsallis per $q<1$.
Sugli indici con $u_i>0$ la funzione $-\sum_ip_i^qu_i$ è strettamente convessa, e il calcolo del caso precedente dà lo stesso risultato con $Z_q=\sum_{u_j>0}u_j^{1/(1-q)}=\sum_j\eq(-x_j/T)$.

Per $q\to1$ si ha $\eq(u)\to e^{u}$ e $\lnq\to\log$.
\end{proof}

Abbiamo verificato numericamente la~\eqref{eq:tsallis-closed} con una minimizzazione diretta sul simplesso per alcuni valori di $(q,T)$.

\begin{corollario}[Limite $T\to0$ della somma di Tsallis]
\label{cor:tsallis-T0}
Siano $x_i>0$.
Per $0<q\le1$ si ha $\lim_{T\to0}F_q(x;T)=\min_ix_i$.
Per $q>1$ si ha
\begin{equation}
\lim_{T\to0}F_q(x;T)=\min_{p\in\Delta_n}\sum_ip_i^{\,q}x_i=\Bigl(\sum_ix_i^{-1/(q-1)}\Bigr)^{-(q-1)}.
\label{eq:power-mean}
\end{equation}
Questo valore è strettamente minore di $\min_ix_i$ se almeno due $x_i$ sono finiti, e tende a $\min_ix_i$ per $q\to1^+$.
\end{corollario}

\begin{proof}
Per $T\to0$ il termine $-TS_q$ scompare, perché $S_q$ è limitata sul simplesso.
Resta il minimo di $U_q(p)=\sum_ip_i^qx_i$.
Per $q<1$ la funzione $U_q$ è concava, quindi il minimo è in un vertice del simplesso e vale $\min_ix_i$.
Per $q=1$ la funzione è lineare e il risultato è lo stesso.
Per $q>1$ la funzione è strettamente convessa.
La stazionarietà dà $p_i\propto x_i^{-1/(q-1)}$, e sostituendo si ottiene la~\eqref{eq:power-mean}.
Ogni termine della somma in~\eqref{eq:power-mean} è positivo, quindi la somma supera $(\min_ix_i)^{-1/(q-1)}$ e il valore è minore di $\min_ix_i$.
Per $q\to1^+$ l'esponente $1/(q-1)$ tende a infinito, e la media di potenze tende al minimo.
\end{proof}

Il corollario dice che la somma di Tsallis deformata con $q>1$ non si dequantizza al semianello tropicale.
Il suo limite $T\to0$ è una media di potenze, non il minimo.
I due limiti iterati danno entrambi il minimo: $\lim_{T\to0}\lim_{q\to1}F_q=\min_ix_i$ e $\lim_{q\to1^+}\lim_{T\to0}F_q=\min_ix_i$.
Il quadrato dei limiti quindi commuta, ma a $q>1$ fisso il limite $T\to0$ ha un vertice diverso dal minimo.
Questo fatto compare nel diagramma della Sezione~\ref{sec:diagrams}.

\subsection{Funzione successore e cumulanti}

MT definiscono la funzione successore $\lambda(x,T)=x\oplus_{S,T}0$ (Sezione~9).
Essa è la trasformata di Legendre di $TS$.
La Proposizione~9.4 di MT afferma che $-\lambda/T$ è la funzione generatrice dei cumulanti dell'energia.
Nel caso di Shannon questa è la relazione usuale fra energia libera e cumulanti.
Nella Sezione~\ref{sec:spy} la useremo nella forma seguente: $G(y)=-y\Phi(y)$ è la funzione generatrice dei cumulanti dell'energia degli stati, per variabile, e la trasformata~\eqref{eq:Phi-y} è la relazione di Gärtner--Ellis fra questa funzione e la complessità~\citep{touchette2009large}.

\section{SP-$y$ passo per passo}
\label{sec:spy}

Questa sezione deriva le equazioni SP-$y$ per il $K$-SAT.
Le equazioni sono di Mézard e Zecchina~\citep{mezard2002random} e di Battaglia, Kolář e Zecchina~\citep{battaglia2004minimizing}.
Le scriviamo con l'energia $1$ per clausola violata.
Le convenzioni sull'energia e sulla normalizzazione delle survey cambiano fra quei lavori, e nel confronto il parametro $y$ può differire per un fattore $2$.
Il contributo di questa sezione è la lettura in termini di semianelli, a due livelli distinti.

\subsection{Livello dei cluster: il potenziale come somma di Shannon}

Sia $\{E_\alpha\}$ l'insieme delle energie degli stati.
Per la Proposizione~\ref{prop:shannon} con $T=1/y$ si ha
\begin{equation}
N\,\Phi_N(y)\equiv-\frac1y\log\sum_\alpha e^{-yE_\alpha}=\bigoplus_{\Sh,1/y}\,\{E_\alpha\}_\alpha
=\min_{\rho}\Bigl[\sum_\alpha\rho_\alpha E_\alpha-\frac1y\,H(\rho)\Bigr],
\label{eq:Phi-shannon}
\end{equation}
dove $\rho$ è una distribuzione sugli stati.
Il minimo è la distribuzione di Gibbs $\rho_\alpha^*\propto e^{-yE_\alpha}$.
Nel limite $N\to\infty$ la distribuzione $\rho^*$ si concentra sugli stati con energia vicina a $Ne^*(y)$.
Il loro numero è $e^{N\Sig_e(e^*)}$, quindi $H(\rho^*)=N\Sig_e(e^*)+o(N)$.
La~\eqref{eq:Phi-shannon} diventa $\Phi=e^*-\Sig_e(e^*)/y$, che è la~\eqref{eq:Phi-y}.

\begin{osservazione}
\label{oss:two-levels}
La complessità $\Sig_e(e^*)$ è l'entropia di Shannon, per variabile, della misura di Gibbs sugli stati alla temperatura $T=1/y$.
Il potenziale $\Phi(y)$ è la somma termodinamica di Shannon delle energie degli stati.
Ci sono due dequantizzazioni.
Il limite $y\to\infty$, cioè $T\to0$, dà il minimo tropicale $\min_\alpha E_\alpha$.
Il limite $N\to\infty$, con parametro $1/N$, dà la trasformata di Legendre~\eqref{eq:Phi-y}.
\end{osservazione}

\subsection{Il reweighting di cavità segue dal livello dei cluster}

Consideriamo il grafo di cavità in cui manca la variabile $j$, e aggiungiamo $j$ con le sue clausole.
Ogni stato $\alpha$ del grafo di cavità genera uno stato del grafo allargato con energia $E_\alpha+\Delta E_\alpha$, dove $\Delta E_\alpha\ge0$ è il numero minimo di clausole violate che l'aggiunta introduce.
Allora
\begin{equation}
\sum_{\alpha}e^{-y(E_\alpha+\Delta E_\alpha)}=\Bigl(\sum_\alpha e^{-yE_\alpha}\Bigr)\ \E_{\rho^*}\bigl[e^{-y\Delta E}\bigr].
\label{eq:reweight}
\end{equation}
La distribuzione delle grandezze locali nel grafo allargato è quindi la distribuzione nel grafo di cavità moltiplicata per $e^{-y\Delta E}$ e normalizzata.
Mézard e Zecchina ottengono lo stesso fattore dal punto di vista microcanonico (eq.~(33) di~\citep{mezard2002random}).
Il numero di stati con energia $E-\Delta E$ nel grafo di cavità è $e^{N\Sig_e((E-\Delta E)/N)}\simeq e^{N\Sig_e(E/N)}\,e^{-y\Delta E}$, con $y=\Sig_e'(e)$.

\begin{osservazione}
\label{oss:no-free-T}
Il parametro $y$ dell'aggregazione locale e il parametro $y$ del potenziale dei cluster sono lo stesso numero.
Una temperatura locale scelta in modo indipendente non corrisponde a nessun potenziale 1RSB.
\end{osservazione}

\subsection{Livello locale: warning, costo e settori}

Fissiamo il lato $(j,a)$.
La variabile $j$ riceve i warning $u_{b\to j}$ dalle clausole $b\in\partial j\setminus a$.
Sia $\omega=(u_{b\to j})_{b\in\partial j\setminus a}$ la configurazione dei warning.
Contiamo i warning che spingono verso il valore che soddisfa $a$ e quelli che spingono verso il valore opposto:
\begin{equation}
p(\omega)=\sum_{b\in\partial^s_aj}u_{b\to j},\qquad q(\omega)=\sum_{b\in\partial^u_aj}u_{b\to j}.
\end{equation}

\paragraph{Passo 1: la distribuzione a priori.}
Nel grafo di cavità i warning sono indipendenti, quindi
\begin{equation}
P_0(\omega)=\prod_{b\in\partial j\setminus a}\bigl[(1-\eta_{b\to j})\,\delta_{u_{b\to j},0}+\eta_{b\to j}\,\delta_{u_{b\to j},1}\bigr].
\label{eq:prior}
\end{equation}

\paragraph{Passo 2: il costo.}
Ogni warning dice che la clausola $b$ ha tutte le altre variabili congelate contro di sé.
Se $x_j$ non segue il warning, la clausola $b$ è violata.
La scelta migliore di $x_j$ segue il lato con più warning, e il costo è
\begin{equation}
\Delta E(\omega)=\min\bigl(p(\omega),q(\omega)\bigr).
\label{eq:cost}
\end{equation}
Le clausole che non mandano warning hanno una variabile libera o una variabile che le soddisfa, e non costano nulla.

\paragraph{Passo 3: i settori.}
Il messaggio da $j$ ad $a$ dipende solo dal lato vincente.
Definiamo tre settori:
\begin{equation}
\Omega_u=\{\omega:\ q>p\},\qquad\Omega_s=\{\omega:\ p>q\},\qquad\Omega_0=\{\omega:\ p=q\}.
\end{equation}
Nel settore $u$ la variabile $j$ è congelata contro $a$.
Nel settore $s$ è congelata a favore di $a$.
Nel settore $0$ entrambi i valori hanno lo stesso costo e $j$ non è congelata.

\paragraph{Passo 4: il reweighting.}
Per la~\eqref{eq:reweight} il peso di $\omega$ è $P_0(\omega)\,e^{-y\Delta E(\omega)}$.
Il peso del settore $\sigma\in\{u,s,0\}$ è
\begin{equation}
\widetilde\Pi^\sigma_{j\to a}(y)=\sum_{\omega\in\Omega_\sigma}P_0(\omega)\,e^{-y\min(p(\omega),q(\omega))}.
\label{eq:sector}
\end{equation}

\paragraph{Passo 5: la funzione generatrice.}
Sia $C_{pq}$ la probabilità a priori di avere $p$ warning dal lato $s$ e $q$ dal lato $u$.
Essa è il coefficiente di $z^pw^q$ nel polinomio
\begin{equation}
\mathcal C_{j\to a}(z,w)=\prod_{b\in\partial^s_aj}\bigl(1-\eta_{b\to j}+\eta_{b\to j}\,z\bigr)\prod_{b\in\partial^u_aj}\bigl(1-\eta_{b\to j}+\eta_{b\to j}\,w\bigr).
\label{eq:genfun}
\end{equation}
Nel settore $u$ il minimo è $p$, nel settore $s$ è $q$, nel settore $0$ è $p=q$.
Quindi
\begin{equation}
\widetilde\Pi^u=\sum_{q>p}C_{pq}\,e^{-yp},\qquad
\widetilde\Pi^s=\sum_{p>q}C_{pq}\,e^{-yq},\qquad
\widetilde\Pi^0=\sum_{p}C_{pp}\,e^{-yp}.
\label{eq:sectors-explicit}
\end{equation}
Il calcolo dei coefficienti $C_{pq}$ costa $O(d^2)$ operazioni per lato, dove $d$ è il grado della variabile.

\paragraph{Passo 6: le equazioni di messaggio.}
Normalizzando si ottiene
\begin{equation}
\nu_{j\to a}(y)=\frac{\widetilde\Pi^u_{j\to a}}{\widetilde\Pi^u_{j\to a}+\widetilde\Pi^s_{j\to a}+\widetilde\Pi^0_{j\to a}}.
\label{eq:spy-var}
\end{equation}
La clausola $a$ manda un warning a $i$ se e solo se tutte le altre variabili sono nel settore $u$.
Nella nostra convenzione il costo dei warning è già contato sul lato della variabile, quindi l'equazione di clausola resta un prodotto:
\begin{equation}
\eta_{a\to i}=\prod_{j\in\partial a\setminus i}\nu_{j\to a}(y).
\label{eq:spy-clause}
\end{equation}

\subsection{Lettura come somma termodinamica}

Il logaritmo del peso di settore è una somma termodinamica con entropia relativa.
Sia $P_0(\cdot\,|\,\Omega_\sigma)$ la distribuzione a priori condizionata al settore $\sigma$.
Per la Proposizione~\ref{prop:kl} con $T=1/y$,
\begin{equation}
-\frac1y\log\widetilde\Pi^\sigma_{j\to a}=-\frac1y\log P_0(\Omega_\sigma)+\bigoplus_{\KL(\cdot\|P_0(\cdot|\Omega_\sigma)),\,1/y}\bigl\{\Delta E(\omega)\bigr\}_{\omega\in\Omega_\sigma}.
\label{eq:sector-kl}
\end{equation}
Il primo termine è il costo di informazione del settore.
Il secondo è l'energia libera dei conflitti dentro il settore.
Il riferimento $P_0$ non è uniforme, quindi l'operazione non è commutativa nel senso di MT: ogni configurazione porta il suo peso a priori.
Questa non commutatività non crea problemi, perché la somma avviene su un insieme finito di configurazioni etichettate.

\subsection{I due limiti in $y$}

\begin{proposizione}[SP come limite $y\to\infty$]
\label{prop:y-infty}
Per $y\to\infty$ le~\eqref{eq:sectors-explicit} tendono alle~\eqref{eq:Pi-u}--\eqref{eq:Pi-0}, e le~\eqref{eq:spy-var}--\eqref{eq:spy-clause} tendono alle equazioni SP.
\end{proposizione}

\begin{proof}
Si ha $e^{-y\min(p,q)}\to\mathbf 1[\min(p,q)=0]$.
Nel settore $u$ resta $p=0<q$, con peso $\sum_{q>0}C_{0q}=\prod_{b\in\partial^s_aj}(1-\eta_{b\to j})\,[1-\prod_{b\in\partial^u_aj}(1-\eta_{b\to j})]=\Pi^u$.
Nel settore $s$ resta $q=0<p$, con peso $\Pi^s$.
Nel settore $0$ resta $p=q=0$, con peso $C_{00}=\Pi^0$.
Le configurazioni con $p,q\ge1$ sono le contraddizioni e hanno peso zero.
\end{proof}

Per $y\to0$ il peso $e^{-y\Delta E}$ tende a $1$, e i settori diventano le probabilità a priori $P_0(q>p)$, $P_0(p>q)$ e $P_0(p=q)$.
Il messaggio segue la maggioranza dei warning e le contraddizioni non costano nulla.
Questo limite è formale.
La descrizione 1RSB degli stati è stabile solo sotto una energia di soglia~\citep{montanari2004instability}, quindi il ramo a $y$ piccolo non descrive direttamente stati fisici.

\subsection{Il funzionale di Bethe a $y$ finito}

La stessa logica dei costi dà il funzionale di Bethe per $G(y)=-y\Phi(y)$.
Ogni termine è $\log\E[e^{-y\Delta E}]$ per l'aggiunta di un nodo o di un lato, con le grandezze in ingresso indipendenti.
\begin{itemize}
\item Clausola $a$: il costo è $1$ se tutte le variabili sono nel settore $u$, e zero altrimenti. Quindi $F_a(y)=\log\bigl[1-(1-e^{-y})\prod_{j\in\partial a}\nu_{j\to a}\bigr]$.
\item Sito $i$: con $p$ warning verso $+1$ e $q$ verso $-1$ il costo è $\min(p,q)$. Quindi $F_i(y)=\log\sum_{p,q}C^{(i)}_{pq}\,e^{-y\min(p,q)}$, dove $C^{(i)}_{pq}$ sono i coefficienti di $\prod_{b\in\partial i^+}(1-\eta_{b\to i}+\eta_{b\to i}z)\prod_{b\in\partial i^-}(1-\eta_{b\to i}+\eta_{b\to i}w)$.
\item Lato $(i,a)$: il costo è $1$ se $i$ è nel settore $u$ rispetto ad $a$ e $a$ manda un warning a $i$. Quindi $F_{ia}(y)=\log\bigl[1-(1-e^{-y})\,\nu_{i\to a}\,\eta_{a\to i}\bigr]$.
\end{itemize}
Il funzionale è
\begin{equation}
N\,G(y)\simeq\sum_aF_a(y)+\sum_iF_i(y)-\sum_{(i,a)}F_{ia}(y).
\label{eq:bethe-y}
\end{equation}
Per il lato si usa che, se $i$ è nel settore $0$ rispetto ad $a$ e riceve il warning di $a$, il nuovo costo è $\min(p+1,p)=p$ e non cambia.
Per $y\to\infty$ si ha $1-e^{-y}\to1$ e $\sum_{p,q}C^{(i)}_{pq}e^{-y\min(p,q)}\to1-\pi_i^+\pi_i^-$, e la~\eqref{eq:bethe-y} tende alla~\eqref{eq:bethe-sp}.
Il termine di clausola coincide con quello di Battaglia, Kolář e Zecchina (eq.~(21) di~\citep{battaglia2004minimizing}), che usano la forma con clausole e siti al posto della forma con lati.
L'energia e la complessità seguono dalle~\eqref{eq:legendre-y}: $e^*=-G'(y)$ e $\Sig_e=G+ye^*$.
Nella fase UNSAT si cerca il punto $y^*$ in cui $\Sig_e=0$~\citep{mezard2002random,battaglia2004minimizing}.

\subsection{Varianti e algoritmi sull'asse $y$}

Sull'asse $y$ si trovano tre algoritmi noti.
Il primo è SID, cioè $y=\infty$ senza backtracking.
Il secondo è BSP, cioè $y=\infty$ con il backtracking governato da $r$ e dal bias~\eqref{eq:bias}.
Il terzo è SP-Y, cioè $y$ finito con il backtracking governato da $r$ e dall'indice~\eqref{eq:bkz-index}.
Battaglia, Kolář e Zecchina propongono anche di aggiornare $y$ durante la decimazione per restare vicino al massimo $y^*$ di $\Phi(y)$.
Nel nostro codice la funzione \texttt{thermo\_star} calcola le somme di settore~\eqref{eq:sectors-explicit} con $T=1/y$.
Con i parametri $\gamma=0$ e $m=0$ la variante chiamata ThermoSP è quindi SP-$y$, e non è un algoritmo nuovo.

\section{Deformazioni sui pesi dei cluster: Rényi, Tsallis e il tilt $\gamma$}
\label{sec:qaxis}

La Sezione~\ref{sec:spy} mostra che la somma di Shannon agisce a due livelli.
Qui deformiamo il livello dei cluster.
Le entropie di Rényi e di Tsallis danno due deformazioni diverse, con conseguenze diverse vicino alla soglia $\alpha_s$.

\subsection{L'ordine di Rényi è il parametro di Parisi}

Nella fase SAT diamo a ogni cluster il peso $p_C=|C|/\mathcal Z$, dove $\mathcal Z=\sum_C|C|$ è il numero di soluzioni.
Questo è il peso del cluster nella misura uniforme sulle soluzioni.

\begin{proposizione}
\label{prop:renyi-m}
L'entropia di Rényi di ordine $m$ dei pesi $\{p_C\}$ è
\begin{equation}
H_m=\frac{1}{1-m}\log\sum_Cp_C^{\,m}=\frac{\log\Xi(m)-m\log\Xi(1)}{1-m},\qquad \Xi(m)=\sum_C|C|^m.
\label{eq:renyi-clusters}
\end{equation}
Per $N\to\infty$ la sua densità è
\begin{equation}
h_m=\lim_{N\to\infty}\frac{H_m}{N}=\frac{\varphi(m)-m\,\varphi(1)}{1-m}.
\label{eq:hm}
\end{equation}
\end{proposizione}

\begin{proof}
Si ha $\sum_Cp_C^m=\mathcal Z^{-m}\sum_C|C|^m=\Xi(m)/\Xi(1)^m$.
Si prende il logaritmo, si divide per $1-m$ e si usa la~\eqref{eq:phi-m}.
\end{proof}

L'ordine della famiglia di Rényi coincide quindi con il parametro di Parisi del limite entropico.
I casi particolari sono i seguenti.
Per $m\to0$ si ha $h_0=\varphi(0)$, l'entropia di Hartley, cioè il logaritmo del numero di cluster: questo è il conteggio di SP.
Per $m\to1$ la regola di de l'Hôpital e la~\eqref{eq:legendre-m} danno $h_1=\varphi(1)-\varphi'(1)=\Sig_s(s^*(1))$, l'entropia di Shannon dei pesi, cioè la complessità dei cluster che dominano la misura uniforme.
Per $m\to\infty$ si ha la min-entropia, $h_\infty=\varphi(1)-s_{\max}$, dove $s_{\max}$ è la massima entropia interna con $\Sig_s\ge0$: essa misura il cluster più grande.

\paragraph{Lettura multifrattale.}
Scriviamo $p_C=e^{-Na}$, con $a=\varphi(1)-s_C\ge0$.
Il numero di cluster con esponente $a$ è $e^{Nf(a)}$ con $f(a)=\Sig_s(\varphi(1)-a)$.
Poniamo $\tau(m)=-\lim N^{-1}\log\sum_Cp_C^m=m\varphi(1)-\varphi(m)$.
Dalla~\eqref{eq:phi-m} si ottiene
\begin{equation}
\tau(m)=\inf_{a}\bigl[m\,a-f(a)\bigr],\qquad h_m=\frac{\tau(m)}{m-1}.
\label{eq:multifractal}
\end{equation}
Questo è il formalismo multifrattale di Halsey e coautori~\citep{halsey1986fractal}, con risoluzione $e^{-N}$.
La funzione $f(a)$ è lo spettro delle singolarità, $\tau(m)$ è la sua trasformata di Legendre, e $h_m$ è la dimensione generalizzata di ordine $m$.
La coppia $(\varphi(m),\Sig_s(s))$ della 1RSB entropica è quindi la stessa coppia di Legendre $(\tau(q),f(\alpha))$ della teoria dei multifrattali.
Il calcolo di $\Sig_s(s)$ tramite $\varphi(m)$ è il metodo di~\citep{mezard2005landscape,krzakala2007clusters}.
MT citano l'uso dell'entropia di Rényi nei multifrattali, ma non collegano l'ordine al parametro di Parisi.

\paragraph{Condensazione e soglia.}
Sia $m^*$ il valore in cui $\Sig_s(s^*(m^*))=0$.
Se $m^*<1$ la misura è condensata, e questo accade per $\alpha_c<\alpha<\alpha_s$~\citep{krzakala2007clusters}.
Per $m\ge m^*$ il supremo in~\eqref{eq:phi-m} è sul bordo $\Sig_s=0$, quindi $\varphi(m)=m\,s_{\max}$ e $\varphi(1)=s_{\max}$.
Sostituendo nella~\eqref{eq:hm} si ottiene
\begin{equation}
h_m=0\qquad\text{per ogni }m\ge m^*,\ m\ne1,
\label{eq:hm-zero}
\end{equation}
e anche $h_1=0$ per continuità.
Nella fase condensata tutte le entropie di Rényi di ordine $m\ge m^*$ sono sotto-estensive: la misura sta su un numero sotto-esponenziale di cluster.
Solo gli ordini $m<m^*$ vedono ancora un numero esponenziale di cluster.
Per $\alpha\to\alpha_s$ la complessità totale tende a zero e $m^*\to0$~\citep{krzakala2007clusters}.
Vicino alla soglia, quindi, lo spettro multifrattale utile si restringe all'intervallo $0\le m<m^*$, che tende al punto di SP.

Questa è la risposta precisa alla domanda sulla statistica multifrattale vicino alla soglia SAT--UNSAT.
Una deformazione con $m>0$ fisso smette di vedere informazione estensiva prima di $\alpha_s$.
Il numero $m^*(\alpha)$ è esso stesso una misura della distanza dalla soglia.
C'è anche un costo pratico.
Con $m>0$ le equazioni di cavità non si chiudono sulle survey: ogni lato porta una distribuzione dei messaggi di BP, e la dimensione del cluster pesa i messaggi~\citep{mezard2005landscape,krzakala2007clusters,mezard2009information}.
Una variante di SP con $m>0$ richiede quindi una popolazione di messaggi per ogni lato.

\subsection{Tsallis sui pesi dei cluster}

Sostituiamo il peso di Boltzmann degli stati con l'esponenziale deformato, come nella Proposizione~\ref{prop:tsallis}:
\begin{equation}
\mathcal Z_q(y)=\sum_\alpha\eq(-yE_\alpha).
\label{eq:Zq}
\end{equation}

\begin{proposizione}[Tre regimi]
\label{prop:tsallis-regimes}
Supponiamo che il numero di stati con energia $E=Ne$ sia $e^{N\Sig_e(e)+o(N)}$, con $\Sig_e$ continua.
\begin{enumerate}
\item Se $q>1$ è fisso, allora $N^{-1}\log\mathcal Z_q(y)\to\sup_e\Sig_e(e)$.
\item Se $q<1$ è fisso, allora $N^{-1}\log\mathcal Z_q(y)\to\Sig_e(0)$, cioè al conteggio degli stati di energia nulla.
\item Se $q-1=\kappa/N$ con $\kappa$ fisso, allora
\begin{equation}
\frac1N\log\mathcal Z_q(y)\to\sup_{e}\Bigl[\Sig_e(e)-g_\kappa(e)\Bigr],\qquad g_\kappa(e)=\frac1\kappa\log(1+\kappa ye),
\label{eq:gen-canonical}
\end{equation}
con il vincolo $e<1/(|\kappa|y)$ se $\kappa<0$.
\end{enumerate}
\end{proposizione}

\begin{proof}
Per $q>1$ si ha $\log\eq(-yNe)=-(q-1)^{-1}\log(1+(q-1)yNe)=O(\log N)$.
Il peso è sotto-esponenziale, e tutte le energie contano allo stesso modo sulla scala $e^{N}$.
Per $q<1$ il peso è zero se $E>1/((1-q)y)$, che è un numero finito.
Contano solo gli stati con $E=O(1)$, cioè $e=0$ sulla scala di $N$.
Per $q-1=\kappa/N$ con $\kappa>0$ si ha $\eq(-yNe)=(1+\kappa ye)^{-N/\kappa}=e^{-Ng_\kappa(e)}$.
Per $\kappa<0$ si ha $\eq(-yNe)=(1-|\kappa|ye)_+^{N/|\kappa|}$, che dà la stessa formula con il taglio.
Il metodo di Laplace conclude.
\end{proof}

I regimi 1 e 2 non danno nulla di nuovo.
Il regime 1 dà $\sup_e\Sig_e(e)=G(0)$, cioè il limite $y\to0$ del potenziale~\eqref{eq:Phi-y}, che le equazioni SP-$y$ non descrivono in modo affidabile (Sezione~\ref{sec:spy}).
Il regime 2 dà il limite $y\to\infty$.
Una deformazione di Tsallis con $q$ fisso sui pesi dei cluster degenera quindi nei due estremi dell'asse $y$.
Il regime 3 è un ensemble canonico generalizzato nel senso di Costeniuc, Ellis, Touchette e Turkington~\citep{costeniuc2005generalized,costeniuc2006generalized}.
In quei lavori il peso di Boltzmann $e^{-\beta H}$ viene moltiplicato per $e^{-g(H)}$ con $g$ continua.

\begin{corollario}[SP-$y$ con un $y$ efficace]
\label{cor:yeff}
Nel regime 3, l'aggiunta di un nodo con spostamento di energia $\Delta E=O(1)$ moltiplica il peso di uno stato per $e^{-y_{\mathrm{eff}}\Delta E}\,(1+O(1/N))$, con
\begin{equation}
y_{\mathrm{eff}}=g_\kappa'(e^*)=\frac{y}{1+\kappa y e^*}.
\label{eq:yeff}
\end{equation}
Le equazioni di messaggio sono quindi le equazioni SP-$y$~\eqref{eq:sectors-explicit}--\eqref{eq:spy-clause} con $y=y_{\mathrm{eff}}$.
L'energia $e^*$ è data dal punto di sella di~\eqref{eq:gen-canonical}, cioè $\Sig_e'(e^*)=y_{\mathrm{eff}}$.
\end{corollario}

\begin{proof}
Si ha $Ng_\kappa((E+\Delta E)/N)-Ng_\kappa(E/N)=g_\kappa'(E/N)\,\Delta E+O(1/N)$, e la misura si concentra su $E/N=e^*$.
Il resto è la derivazione della Sezione~\ref{sec:spy} con $y$ sostituito da $y_{\mathrm{eff}}$.
\end{proof}

Il corollario ha due conseguenze.
Primo, se $\Sig_e$ è concava, il regime 3 non raggiunge punti nuovi.
Ogni punto $e^*$ con $\Sig_e'(e^*)=y_{\mathrm{eff}}$ si raggiunge anche con SP-$y$ al valore $y_{\mathrm{eff}}$, quindi la deformazione è solo una riparametrizzazione dell'asse $y$.
Secondo, il punto $e^*$ è un massimo di $\Sig_e-g_\kappa$ se $\Sig_e''(e^*)<g_\kappa''(e^*)=-\kappa y_{\mathrm{eff}}^2$.
Per $\kappa>0$, cioè $q>1$, la funzione $g_\kappa$ è concava, e il vincolo è più forte di quello di SP-$y$.
Per $\kappa<0$, cioè $q<1$, la funzione $g_\kappa$ è convessa, e il regime 3 raggiunge anche punti in cui $\Sig_e$ è localmente convessa, con $0\le\Sig_e''(e^*)<|\kappa|y_{\mathrm{eff}}^2$.
Questi punti non si raggiungono con nessun valore di $y$.
Questo è il meccanismo dell'ensemble gaussiano di~\citep{costeniuc2006generalized}, con una funzione $g$ diversa.

\paragraph{Uso vicino alla soglia.}
La deformazione con $q<1$ in scala $1/N$ ha anche un taglio di energia $e<1/(|\kappa|y)$.
Essa quindi seleziona una finestra di energie e penalizza di più le energie vicine al taglio.
In un algoritmo la deformazione si realizza con un ciclo esterno su $y_{\mathrm{eff}}$: si risolve SP-$y$ al valore corrente, si calcola $e^*=-G'(y_{\mathrm{eff}})$ dal funzionale~\eqref{eq:bethe-y}, e si aggiorna $y_{\mathrm{eff}}\leftarrow y/(1+\kappa ye^*)$.
Il ciclo è simile all'aggiornamento dinamico di $y$ proposto da Battaglia, Kolář e Zecchina.
Non sappiamo se la complessità energetica $\Sig_e(e)$ del $K$-SAT casuale abbia tratti non concavi vicino ad $\alpha_s$.
Solo in quel caso la deformazione di Tsallis in scala $1/N$ darebbe stati nuovi.
Il nostro codice non contiene ancora questo parametro.

\begin{osservazione}[Tsallis a livello locale]
Si possono sostituire i pesi $e^{-y\min(p,q)}$ nelle~\eqref{eq:sectors-explicit} con $\eq(-y\min(p,q))$ a $q$ fisso.
Il Corollario~\ref{cor:yeff} mostra che questa sostituzione non segue da nessun potenziale sui cluster della forma~\eqref{eq:Zq}.
Sulla scala esponenziale il reweighting di cavità è sempre esponenziale nello spostamento di energia.
Questa variante è quindi un'euristica e non un'equazione 1RSB.
\end{osservazione}

\subsection{Il tilt verso gli stati $*$ nel modello di MMW}

Gli autori di BSP suggeriscono che una misura concentrata sui cluster non congelati potrebbe ridurre il backtracking~\citep{marino2016bsp}.
Il modello~\eqref{eq:mmw} offre un modo esatto per scrivere una misura di questo tipo.
Si prende $\omega_o=0$ e $\omega_*=e^{\gamma}$.
Allora il peso di un assegnamento parziale senza variabili libere è $e^{\gamma n_*}$, e per $\gamma>0$ la misura preferisce gli assegnamenti con molte variabili $*$, cioè con poche variabili congelate.

\begin{proposizione}
\label{prop:gamma}
Sia $\omega_o=0$.
La riduzione di BP sul modello~\eqref{eq:mmw} a una sola survey per lato vale per messaggi generici se e solo se $\omega_*=1$.
\end{proposizione}

\begin{proof}
Usiamo le equazioni di BP sul modello esteso nella Sezione~3 di MMW~\citep{maneva2007survey}.
I messaggi da clausola a variabile hanno le componenti $M^s$, $M^u$ e $M^*$, e i messaggi da variabile a clausola hanno le componenti $R^s$, $R^u$ e $R^*$.
Le componenti $R^s$ e $R^*$ del messaggio da $i$ ad $a$ valgono
\begin{align*}
R^s_{i\to a}&=\prod_{b\in\uu}M^u_{b\to i}\prod_{b\in\us}\bigl(M^s_{b\to i}+M^*_{b\to i}\bigr),\\
R^*_{i\to a}&=\prod_{b\in\uu}M^u_{b\to i}\Bigl[\prod_{b\in\us}\bigl(M^s_{b\to i}+M^*_{b\to i}\bigr)-(1-\omega_o)\prod_{b\in\us}M^*_{b\to i}\Bigr]+\omega_*\prod_{b\in\partial i\setminus a}M^*_{b\to i}.
\end{align*}
Supponiamo che i messaggi in ingresso soddisfino $M^u=M^*$.
Allora
\begin{equation*}
R^s_{i\to a}-R^*_{i\to a}=(1-\omega_o-\omega_*)\prod_{b\in\partial i\setminus a}M^*_{b\to i}.
\end{equation*}
Le equazioni di clausola di MMW danno
\begin{equation*}
M^u_{a\to i}-M^*_{a\to i}=\sum_{k\in\partial a\setminus i}\bigl(R^s_{k\to a}-R^*_{k\to a}\bigr)\prod_{j\in\partial a\setminus\{i,k\}}R^u_{j\to a}.
\end{equation*}
Quindi l'invariante $M^u=M^*$ si conserva se $\omega_o+\omega_*=1$, e per messaggi generici si rompe altrimenti.
Questa è la condizione~(a) del Teorema~3 di MMW.
Con $\omega_o=0$ la condizione diventa $\omega_*=1$.
\end{proof}

Con $\gamma\ne0$ servono quindi i messaggi completi di MMW, con due numeri liberi per lato invece di uno.
Il calcolo è BP su un modello grafico esplicito, quindi è ben definito.
Una chiusura su un solo numero, che moltiplica soltanto il peso $\Pi^0$ per $e^{\gamma}$, è un'approssimazione e non coincide con BP sul modello di MMW.
Il nostro codice usa questa chiusura approssimata.
Il valore algoritmico del tilt $\gamma$ è una domanda aperta.

\section{Risposta locale della complessità}
\label{sec:response}

Questa sezione studia la variazione del funzionale~\eqref{eq:bethe-sp} quando si fissa o si libera una variabile.
L'identità della Proposizione~\ref{prop:local}, nel caso del valore più probabile, è l'affermazione di BSP sul bias~\eqref{eq:bias}.
La forma con il segno, il resto di secondo ordine e il confronto con l'indice~\eqref{eq:bkz-index} sono nostri.

\subsection{Il funzionale con una variabile bloccata}

Sia $F$ la formula corrente e sia $\eta^\star$ un punto fisso di SP su $F$.
Scegliamo una variabile $k$ e un valore $s\in\{-1,+1\}$.
La formula ridotta $F|_{x_k=s}$ si ottiene eliminando $k$, eliminando le clausole in $\partial k^{s}$, che sono soddisfatte, e accorciando le clausole in $\partial k^{-s}$.

Scriviamo il funzionale di $F|_{x_k=s}$ sul grafo di $F$.
Definiamo $\Sig^{(k,s)}[\eta]$ come il funzionale~\eqref{eq:bethe-sp} con tre modifiche.
\begin{enumerate}
\item Per $a\in\partial k^{-s}$ il messaggio $\nu_{k\to a}$ vale $1$: la variabile $k$ non soddisfa $a$.
\item Per $a\in\partial k^{s}$ il messaggio $\nu_{k\to a}$ vale $0$: la variabile $k$ soddisfa $a$, e $a$ non manda più warning.
\item Il termine di sito di $k$ diventa $F_k^{(s)}=\log\prod_{b\in\partial k^{-s}}(1-\eta_{b\to k})=\log(1-\pi_k^{-s})$.
\end{enumerate}

\begin{lemma}
\label{lem:clamped}
Se $\eta_{a\to j}=0$ per ogni $a\in\partial k^{s}$ e $j\in\partial a\setminus k$, allora $\Sig^{(k,s)}[\eta]$ è il funzionale~\eqref{eq:bethe-sp} della formula ridotta $F|_{x_k=s}$.
Inoltre $\Sig^{(k,s)}$ non dipende dalle survey $\eta_{b\to k}$.
\end{lemma}

\begin{proof}
Per $a\in\partial k^{s}$ si ha $F_a=\log(1-0)=0$ e $F_{ka}=0$.
Se inoltre $\eta_{a\to j}=0$, anche $F_{ja}=0$, e la clausola $a$ non contribuisce ai siti $j$: è come se $a$ non ci fosse.
Per $a\in\partial k^{-s}$ si ha $F_a=\log(1-\prod_{j\in\partial a\setminus k}\nu_{j\to a})$, che è il termine della clausola accorciata.
I lati $(k,a)$ con $a\in\partial k^{-s}$ danno $-\sum_{a\in\partial k^{-s}}\log(1-\eta_{a\to k})=-F_k^{(s)}$, che cancella il termine di sito.
Restano i termini della formula ridotta.
Le survey $\eta_{b\to k}$ compaiono solo in $F_k^{(s)}$ e nei lati $(k,b)$, e questi termini si cancellano o valgono zero.
\end{proof}

\subsection{L'identità locale}

\begin{proposizione}[Identità locale]
\label{prop:local}
Siano $\pi_k^{\pm}$ e $w_k^{\pm}$ le grandezze~\eqref{eq:w} calcolate con le survey $\eta^\star_{b\to k}$.
Allora
\begin{equation}
\Sig^{(k,s)}[\eta^\star]-\Sig[\eta^\star]=\log\frac{1-\pi_k^{-s}}{1-\pi_k^+\pi_k^-}=\log\bigl(1-w_k^{-s}\bigr).
\label{eq:local-identity}
\end{equation}
\end{proposizione}

\begin{proof}
I due funzionali differiscono in due punti.
Il primo è il termine di sito di $k$, che passa da $\log(1-\pi_k^+\pi_k^-)$ a $\log(1-\pi_k^{-s})$.
Il secondo sono i messaggi $\nu_{k\to a}$, che entrano solo in $F_a-F_{ka}$.
Nel punto fisso si ha $\eta^\star_{a\to k}=\prod_{j\in\partial a\setminus k}\nu_{j\to a}$, quindi
\begin{equation*}
F_a-F_{ka}=\log\bigl(1-\nu_{k\to a}\eta^\star_{a\to k}\bigr)-\log\bigl(1-\nu_{k\to a}\eta^\star_{a\to k}\bigr)=0
\end{equation*}
per ogni valore di $\nu_{k\to a}$.
Questa è l'identità di Bethe $F_a-F_{ka}=\log c_{a\to k}$, dove $c_{a\to k}$ non dipende da $\nu_{k\to a}$.
Resta solo la variazione del termine di sito.
L'ultima uguaglianza segue da $1-w_k^{-s}=[1-\pi_k^+\pi_k^--\pi_k^{-s}(1-\pi_k^{s})]/(1-\pi_k^+\pi_k^-)=(1-\pi_k^{-s})/(1-\pi_k^+\pi_k^-)$.
\end{proof}

Sia ora $\eta^{(s)}$ il punto fisso di SP sulla formula ridotta, esteso con $\eta_{a\to j}=0$ per $a\in\partial k^{s}$.
La variazione vera della complessità è
\begin{equation}
\Delta\Sig_{k,s}=\Sig^{(k,s)}[\eta^{(s)}]-\Sig[\eta^\star]=\log\bigl(1-w_k^{-s}\bigr)+R_{k,s},
\qquad R_{k,s}=\Sig^{(k,s)}[\eta^{(s)}]-\Sig^{(k,s)}[\eta^\star].
\label{eq:true-change}
\end{equation}

\begin{proposizione}[Il resto è di secondo ordine]
\label{prop:remainder}
Il funzionale $\Sig^{(k,s)}$ è stazionario in $\eta^{(s)}$.
Quindi, con $\delta=\eta^\star-\eta^{(s)}$ e $\mathsf H$ la matrice hessiana di $\Sig^{(k,s)}$ in $\eta^{(s)}$,
\begin{equation}
R_{k,s}=-\tfrac12\,\delta^{\top}\mathsf H\,\delta+O(\|\delta\|^3).
\label{eq:remainder}
\end{equation}
\end{proposizione}

\begin{proof}
Sui lati della formula ridotta il Lemma~\ref{lem:clamped} e il Lemma~\ref{lem:stationary} danno la stazionarietà.
I termini in più, cioè i lati $(j,a)$ con $a\in\partial k^{s}$, hanno derivata nulla rispetto a queste survey, perché contengono il fattore $\eta_{a\to j}=0$.
Restano le survey $\eta_{a\to j}$ con $a\in\partial k^{s}$.
La dimostrazione del Lemma~\ref{lem:stationary} si ripete con due differenze.
La cancellazione fra termine di sito e termine di lato vale per ogni valore delle survey.
Il termine $F_a$ vale zero, quindi il contributo tramite $\nu_{j\to a}$ è $\eta_{a\to j}/(1-\nu_{j\to a}\eta_{a\to j})$, che è zero in $\eta_{a\to j}=0$.
Lo sviluppo di Taylor di $\Sig^{(k,s)}[\eta^\star]$ intorno a $\eta^{(s)}$ non ha quindi termine lineare.
\end{proof}

La proposizione ha una conseguenza diretta per le stime di risposta lineare.
Una correzione della forma $g^{\top}(\mathsf I-\mathsf J)^{-1}b$, dove $g$ è il gradiente del funzionale, $\mathsf J$ è lo jacobiano delle equazioni SP e $b$ è la perturbazione locale, vale zero, perché $g=0$ nel punto fisso.
La prima correzione all'identità~\eqref{eq:local-identity} è di secondo ordine nella variazione dei messaggi.
Il Lemma~\ref{lem:stationary} mostra anche che la forma compatta~\eqref{eq:sigma-bsp} non va usata per questi sviluppi, perché non è stazionaria.

\begin{corollario}[Alberi]
\label{cor:tree}
Se il factor graph di $F$ è un albero, allora $R_{k,s}=0$.
\end{corollario}

\begin{proof}[Idea della dimostrazione]
Su un albero BP è esatta, e il funzionale nel punto fisso è il logaritmo del numero di assegnamenti parziali validi senza variabili libere del modello di MMW con $(\omega_o,\omega_*)=(0,1)$~\citep{braunstein2004surveybp,maneva2007survey}.
Il numero si fattorizza sui sottoalberi appesi a $k$.
Per l'invariante $M^u=M^*$ il peso di un sottoalbero che non manda warning a $k$ non dipende dallo stato di $k$.
Con $k$ libera, la somma sugli stati di $k$ dà il fattore $1-\pi_k^+\pi_k^-$.
Con $k$ bloccata a $s$, le clausole in $\partial k^{-s}$ non possono mandare warning, e il fattore è $1-\pi_k^{-s}$.
Il rapporto è $1-w_k^{-s}$, e la~\eqref{eq:true-change} dà $R_{k,s}=0$.
\end{proof}

Su un albero quindi $R_{k,s}=0$ anche se $\delta\ne0$: le survey che escono da $k$ cambiano, ma il funzionale non cambia.
Su un grafo con cicli ci aspettiamo che $R_{k,s}$ dipenda soprattutto dalle variazioni che tornano nell'intorno di $k$ lungo i cicli.
Questa è un'ipotesi, non un risultato.
Abbiamo verificato numericamente il Lemma~\ref{lem:stationary} e la~\eqref{eq:true-change} su una formula 3-SAT casuale con $N=600$ e $\alpha=4{,}15$.
Il gradiente direzionale del funzionale~\eqref{eq:bethe-sp} è nullo entro l'errore numerico, mentre quello della forma~\eqref{eq:sigma-bsp} non lo è.
Il resto $R_{k,s}$ è piccolo quando $s$ è il valore più probabile, e nel verso opposto può essere dello stesso ordine del termine principale.

\subsection{Conseguenze per la decimazione e il backtracking}

\paragraph{Fissaggio.}
Se $s$ è il valore più probabile, allora $w_k^{-s}=\min(w_k^+,w_k^-)$ e la~\eqref{eq:local-identity} dà $\log b_k$.
Questa è l'affermazione di BSP: il numero di cluster viene moltiplicato per $b_k$~\citep{marino2016bsp}.
La regola di SID e di BSP, che fissa le variabili con $b_k$ più grande, è quindi la regola che minimizza la perdita di complessità stimata.

\paragraph{Degenerazione.}
Si ha $\log(1-w_k^{-s})=0$ quando $w_k^{-s}=0$.
Questo accade in due casi molto diversi: $k$ è congelata a $s$ in tutti i cluster ($w_k^{s}=1$), oppure $k$ non è congelata in nessun cluster ($w_k^0=1$).
In entrambi i casi $b_k=1$.
Il criterio basato sulla sola complessità non distingue i due casi, e serve un secondo criterio per ordinarli, per esempio la polarizzazione $|w_k^+-w_k^-|$.

\paragraph{Rilascio.}
Sia $k$ una variabile già fissata al valore $s_k$ nella formula corrente $F$.
Sia $F'$ la formula in cui $k$ è libera.
La~\eqref{eq:true-change}, applicata a $F'$, dà
\begin{equation}
\Sig(F')-\Sig(F)=-\log\bigl(1-w_k^{-s_k}\bigr)-R_{k,s_k},
\label{eq:release}
\end{equation}
dove $w_k^{\pm}$ sono calcolati con le survey che entrano in $k$.
BSP calcola queste survey con le variabili fissate e le usa per il bias~\citep{marino2016bsp}.
Il guadagno stimato del rilascio è quindi $-\log I_k$ con
\begin{equation}
I_k=1-w_k^{-s_k}.
\label{eq:Ik}
\end{equation}
Il bias di BSP coincide con $I_k$ solo se $s_k$ è ancora il valore più probabile.
Se il campo si è rovesciato, cioè $w_k^{-s_k}>w_k^{s_k}$, allora $b_k=1-w_k^{s_k}>I_k$.
Per esempio, con $w_k^{-s_k}=0{,}9$ e $w_k^{s_k}=0{,}05$ si ha $I_k=0{,}1$ e $b_k=0{,}95$.
La variabile ha un guadagno stimato di $\log10$, ma il bias di BSP la classifica come sicura.
L'indice di Battaglia, Kolář e Zecchina~\eqref{eq:bkz-index} vale $w_k^{-s_k}-w_k^{s_k}$ a $y=\infty$ e tiene conto del segno.
Esso cresce con $w_k^{-s_k}$ come $-\log I_k$, ma sottrae anche $w_k^{s_k}$, quindi i due ordinamenti in generale sono diversi.
Il punto nuovo di $I_k$ non è il segno, che è già in~\citep{battaglia2004minimizing}, ma il legame esatto con la complessità dato dalla~\eqref{eq:local-identity}.

\paragraph{Una famiglia di politiche.}
Le regole di scelta si possono scrivere come distribuzioni di Gibbs sui punteggi.
Per il rilascio, per esempio, si sceglie la variabile $k$ con probabilità proporzionale a $\exp[-\log I_k/\Tact]$, dove $\Tact$ è una temperatura di azione.
Per $\Tact\to0$ si ottiene la scelta deterministica del punteggio migliore.
Questa temperatura non ha legame con $y$: essa regola la politica di ricerca, non la misura sui cluster.

\section{Un look-ahead soft a due passi}
\label{sec:softq}

Questa sezione propone una variante di BSP.
È la sola parte algoritmica del lavoro.
La variante usa le equazioni di Bellman soft del Soft Actor-Critic (SAC) di Haarnoja e coautori~\citep{haarnoja2018soft} e le combina con l'identità locale della Proposizione~\ref{prop:local}.
Il ponte con il resto del lavoro è la somma termodinamica di Shannon: SAC la applica alle azioni della ricerca, come SP la applica ai warning e il potenziale 1RSB agli stati.
Gli enunciati algebrici di questa sezione sono verificati in Lean~4 con Mathlib, nel file \texttt{lean/SoftQ.lean} del codice.
Fra parentesi quadre diamo il nome del teorema Lean corrispondente.

\subsection{Il problema di controllo}

Lo stato è la formula residua $F$, con il suo punto fisso di SP.
Un'azione è una coppia $(k,s)$: si fissa la variabile libera $k$ al valore $s$.
La transizione è deterministica e porta nella formula ridotta $F^{(k,s)}=F|_{x_k=s}$, dopo la propagazione unitaria.
La ricompensa $r(F,(k,s))$ sarà definita nella~\eqref{eq:tsallis-reward}.
Una policy $\pi(\cdot\mid F)$ è una distribuzione sulle azioni.
SAC massimizza la somma delle ricompense più l'entropia della policy, pesata con una temperatura $\alpha>0$.
Con transizioni deterministiche le equazioni di Bellman soft sono
\begin{equation}
Q_\alpha(F,a)=r(F,a)+V_\alpha(F^{a}),\qquad
V_\alpha(F)=\alpha\log\sum_a e^{Q_\alpha(F,a)/\alpha},\qquad
\pi_\alpha(a\mid F)=e^{[Q_\alpha(F,a)-V_\alpha(F)]/\alpha}.
\label{eq:soft-bellman}
\end{equation}
La policy di SAC è quindi la distribuzione di Gibbs dei valori $Q_\alpha$, e vale l'identità $Q_\alpha-V_\alpha=\alpha\log\pi_\alpha$.

\begin{proposizione}[Il valore soft è una somma termodinamica]
\label{prop:sac-shannon}
Sia $Q_a$ un numero reale per ogni azione $a$ di un insieme finito, e sia $\alpha>0$.
Per ogni distribuzione $p$ sulle azioni vale
\begin{equation}
\sum_ap_aQ_a+\alpha H(p)\le\alpha\log\sum_ae^{Q_a/\alpha},
\label{eq:sac-bound}
\end{equation}
e l'uguaglianza vale per $p_a\propto e^{Q_a/\alpha}$.
[\texttt{sac\_objective\_le\_soft}, \texttt{sac\_objective\_gibbs}]
\end{proposizione}

\begin{proof}
È la Proposizione~\ref{prop:shannon} con $x_a=-Q_a$ e $T=\alpha$, letta con il segno opposto.
\end{proof}

Il valore $V_\alpha$ è quindi la somma di Maslov $\oplus_\alpha$ della Sezione~\ref{sec:notation} applicata ai valori delle azioni, cioè $V_\alpha=-\bigoplus_{\Sh,\alpha}\{-Q_a\}_a$.
Per $\alpha\to0$ esso tende a $\max_aQ_a$, e la policy diventa la scelta deterministica della migliore azione.
Nel lavoro la somma di Shannon compare quindi su tre livelli.
Il primo è l'aggregazione dei warning (Sezione~\ref{sec:spy}).
Il secondo è il potenziale dei cluster, dove la temperatura è $1/y$.
Il terzo è la scelta delle mosse della ricerca, dove la temperatura è $\alpha$.
Per l'Osservazione~\ref{oss:no-free-T} i primi due livelli hanno la stessa temperatura.
Il terzo livello non ha questo vincolo: $\alpha$ è la temperatura di azione $\Tact$ della Sezione~\ref{sec:response}, e non corrisponde a nessun potenziale 1RSB.

\subsection{Tre ostacoli a una traduzione letterale}

\begin{proposizione}[La complessità telescopa]
\label{prop:telescoping}
Sia $F_0,F_1,\dots,F_n$ una traiettoria della decimazione e sia $r_t=\Sig(F_{t+1})-\Sig(F_t)$.
Allora $\sum_{t<n}r_t=\Sig(F_n)-\Sig(F_0)$.
Due traiettorie con gli stessi estremi hanno la stessa somma.
[\texttt{telescoping}, \texttt{same\_return}]
\end{proposizione}

Una ricompensa di questa forma è una differenza di potenziali.
Ng, Harada e Russell~\citep{ng1999policy} mostrano che una ricompensa di questo tipo non cambia la policy ottima.
La richiesta di una discesa lenta della complessità è quindi una richiesta sulla forma della traiettoria, e non sulla perdita totale.

\begin{osservazione}[Il look-ahead a un passo]
\label{oss:one-step}
Si fissi $x_k=s$ e si riconverga SP.
Per le Proposizioni~\ref{prop:local} e~\ref{prop:remainder} la variazione misurata della complessità è $\log(1-w_k^{-s})+R_{k,s}$, con $R_{k,s}$ di secondo ordine.
Un look-ahead a un passo sulla complessità ricalcola quindi il bias $b_k$, a meno di un termine di secondo ordine, al prezzo di una riconvergenza completa.
\end{osservazione}

Il secondo passo letterale di SAC è $V_\alpha(F^{(k,s)})=\alpha\log\sum_{l,t}e^{r_l^t/\alpha}$, dove la somma corre sulle mosse $(l,t)$ disponibili dopo la prima.
La proposizione seguente mostra che anche questo valore porta poca informazione.

\begin{proposizione}[Il secondo passo letterale degenera]
\label{prop:degenerate}
\begin{enumerate}
\item Siano $b_l\in[0,1]$ per ogni $l$ di un insieme finito, e sia $b_{l_0}=1$ per un $l_0$. Allora $\max_l\log b_l=0$. [\texttt{hard\_two\_step\_degenerate}]
\item Siano $Z>0$, $Z+\delta>0$ e $\alpha\ge0$. Allora
\begin{equation}
\frac{\alpha\delta}{Z+\delta}\le\alpha\log(Z+\delta)-\alpha\log Z\le\frac{\alpha\delta}{Z}.
\label{eq:dilution}
\end{equation}
[\texttt{dilution}]
\end{enumerate}
\end{proposizione}

\begin{proof}
Il primo punto segue da $\log b_l\le0$ e $\log b_{l_0}=0$.
Il secondo segue da $1-1/x\le\log x\le x-1$ con $x=(Z+\delta)/Z$.
\end{proof}

Nel limite $\alpha\to0$ il secondo passo vale $\max_{l}\log b_l$.
Durante quasi tutta la decimazione esiste una variabile con $b_l=1$, congelata o libera in tutti i cluster, e il primo punto dà zero per ogni prima mossa.
Per $\alpha>0$ la somma $Z=\sum_{l,t}e^{r_l^t/\alpha}$ ha $2N_t$ termini.
La prima mossa cambia solo i termini delle variabili vicine a $k$, quindi $\delta=O(1)$.
Se una frazione finita dei termini è di ordine uno, allora $Z=\Theta(N_t)$, e la~\eqref{eq:dilution} dà una differenza di ordine $1/N_t$ fra due prime mosse.

\subsection{La ricompensa di Tsallis}

Sia $\lnq x=(x^{1-q}-1)/(1-q)$ il logaritmo di Tsallis della Proposizione~\ref{prop:tsallis}.
Definiamo la ricompensa della mossa $(k,s)$ come
\begin{equation}
r_k^s=\lnq\bigl(1-w_k^{-s}\bigr).
\label{eq:tsallis-reward}
\end{equation}
Per $q=1$ essa è la variazione della complessità data dall'identità locale~\eqref{eq:local-identity}.

\begin{proposizione}[Proprietà della ricompensa]
\label{prop:tsallis-reward}
Siano $x,y>0$ e $q\ne1$.
\begin{enumerate}
\item $\lnq(xy)=\lnq x+\lnq y+(1-q)\lnq x\,\lnq y$. [\texttt{lnq\_mul}]
\item Se $q>1$ e $x,y<1$, allora $\lnq(xy)<\lnq x+\lnq y$. [\texttt{lnq\_split}]
\item Se $q>1$, la funzione $\lnq$ è strettamente crescente. [\texttt{lnq\_strictMono}]
\item Per $q\to1$ si ha $\lnq x\to\log x$. [\texttt{lnq\_tendsto\_log}]
\end{enumerate}
\end{proposizione}

\begin{proof}
Il primo punto segue da $(xy)^{1-q}=x^{1-q}y^{1-q}$.
Per il secondo, con $q>1$ e $x<1$ si ha $x^{1-q}>1$ e quindi $\lnq x<0$; lo stesso vale per $y$, e il termine $(1-q)\lnq x\,\lnq y$ è negativo.
Il terzo segue da $x^{1-q}$ decrescente e da $1-q<0$.
Il quarto è la derivata di $t\mapsto x^t$ in $t=0$, con $t=1-q$.
\end{proof}

Il primo punto dice che per $q\ne1$ la somma delle ricompense non telescopa.
Il secondo dice che per $q>1$ una caduta di un fattore $xy$ in un solo passo costa più di due cadute di fattori $x$ e $y$.
La ricompensa premia quindi una discesa lenta e regolare.
Il terzo dice che, su un solo passo, la ricompensa ordina le mosse come il bias $b_k$.
La deformazione di Tsallis agisce qui sulla ricompensa della ricerca.
Nella Sezione~\ref{sec:qaxis} essa agiva invece sui pesi dei cluster, e le due deformazioni non vanno confuse.

\subsection{Il valore a campi congelati}

Dopo la prima mossa approssimiamo il futuro con i campi congelati.
Questo significa che la ricompensa di una mossa futura $(l,t)$ è $r_l^t$ calcolata con le grandezze $w_l^\pm$ di $F^{(k,s)}$, e non dipende dalle mosse fatte prima di essa.

\begin{proposizione}[Fattorizzazione]
\label{prop:frozen}
Sia $L$ un insieme finito di variabili, e siano $r_l^\pm$ numeri reali.
\begin{enumerate}
\item Il guadagno di un assegnamento completo $t\in\{\pm1\}^L$ è $\sum_lr_l^{t_l}$, per ogni ordine in cui le variabili vengono fissate. [\texttt{frozen\_order\_free}]
\item Per ogni $\alpha>0$ vale
\begin{equation}
\alpha\log\sum_{t\in\{\pm1\}^L}\exp\Bigl(\frac1\alpha\sum_{l\in L}r_l^{t_l}\Bigr)=\sum_{l\in L}v_\alpha(l),\qquad
v_\alpha(l)=\alpha\log\bigl(e^{r_l^+/\alpha}+e^{r_l^-/\alpha}\bigr).
\label{eq:frozen-value}
\end{equation}
[\texttt{frozen\_soft\_value}]
\end{enumerate}
\end{proposizione}

\begin{proof}
Il primo punto è l'invarianza di una somma finita per permutazione.
Il secondo segue da $\sum_t\prod_le^{r_l^{t_l}/\alpha}=\prod_l\sum_{t_l}e^{r_l^{t_l}/\alpha}$ e dal logaritmo di un prodotto.
\end{proof}

Il valore del futuro a campi congelati è quindi
\begin{equation}
V_\alpha(F)=\sum_{l\ \text{libera in}\ F}v_\alpha(l).
\label{eq:frozen-V}
\end{equation}
Abbiamo tolto il fattore $N_t!$ degli ordinamenti, perché ordini diversi portano allo stesso assegnamento.
Il valore~\eqref{eq:frozen-V} è estensivo, a differenza del secondo passo letterale.

\begin{proposizione}[Limite duro del valore di sito]
\label{prop:site-limit}
Per $\alpha>0$ vale $\max(r_l^+,r_l^-)\le v_\alpha(l)\le\max(r_l^+,r_l^-)+\alpha\log2$.
Quindi $v_\alpha(l)\to\max(r_l^+,r_l^-)=\lnq b_l$ per $\alpha\to0^+$, se $q\ge1$.
[\texttt{soft2\_ge\_max}, \texttt{soft2\_le\_max\_add}, \texttt{soft2\_tendsto\_max}]
\end{proposizione}

\begin{proof}
La somma $e^{r_l^+/\alpha}+e^{r_l^-/\alpha}$ sta fra $e^{m/\alpha}$ e $2e^{m/\alpha}$, con $m=\max(r_l^+,r_l^-)$.
L'uguaglianza $m=\lnq b_l$ segue dalla monotonia di $\lnq$ e da $b_l=1-\min(w_l^+,w_l^-)$.
\end{proof}

\subsection{Il punteggio}

Il valore soft della mossa $(k,s)$ è $Q_\alpha(F,(k,s))=r_k^s+V_\alpha(F^{(k,s)})$.
Siano $L$ e $L'$ gli insiemi delle variabili libere in $F$ e in $F^{(k,s)}$.
L'insieme $L'$ non contiene $k$ e le variabili fissate dalla propagazione unitaria.

\begin{proposizione}[Decomposizione del punteggio]
\label{prop:score}
Vale
\begin{equation}
Q_\alpha(F,(k,s))-V_\alpha(F)=\alpha\log\pi_\alpha(s\mid k)+D_\alpha(k,s),
\label{eq:score}
\end{equation}
con
\begin{equation}
\pi_\alpha(s\mid k)=\frac{e^{r_k^s/\alpha}}{e^{r_k^+/\alpha}+e^{r_k^-/\alpha}},\qquad
D_\alpha(k,s)=\sum_{l\in L'}v_\alpha\bigl(l;F^{(k,s)}\bigr)-\sum_{l\in L\setminus k}v_\alpha(l;F).
\label{eq:D}
\end{equation}
[\texttt{score\_decomposition}, \texttt{sac\_identity}]
\end{proposizione}

\begin{proof}
Si separa il termine di $k$ in $V_\alpha(F)$, e si usa $\alpha\log\pi_\alpha(s\mid k)=r_k^s-v_\alpha(k)$.
\end{proof}

Il primo termine è l'identità di SAC applicata alla scelta del segno di $k$.
Il secondo termine è il vero look-ahead: misura quanto la mossa cambia il valore soft delle altre variabili.

\begin{corollario}[Variabili congelate e variabili libere]
\label{cor:frozen-free}
Se $k$ non è congelata in nessun cluster, cioè $w_k^+=w_k^-=0$, allora $v_\alpha(k)=\alpha\log2$ e $\alpha\log\pi_\alpha(s\mid k)=-\alpha\log2$.
Se $w_k^{-s}=0$, allora $r_k^s=0$ e $0\le v_\alpha(k)\le\alpha e^{r_k^{-s}/\alpha}$.
Se inoltre $w_k^s\to1$, cioè $k$ è congelata al valore $s$ in quasi tutti i cluster, allora $r_k^{-s}\to-\infty$ e $\alpha\log\pi_\alpha(s\mid k)\to0$.
[\texttt{soft2\_unfrozen}, \texttt{soft2\_frozen}]
\end{corollario}

Le due variabili del corollario hanno entrambe $b_k=1$, e la Sezione~\ref{sec:response} ha mostrato che il criterio della complessità non le distingue.
Il punteggio soft le distingue: per $\alpha>0$ esso fissa prima le variabili congelate e lascia libere le variabili libere in tutti i cluster.
Questo è lo stesso ordine del secondo criterio di BSP, la polarizzazione $|w_k^+-w_k^-|$.

\begin{osservazione}[Il termine a un passo si cancella]
\label{oss:cancel}
Per $\alpha\to0$ e $s$ il valore più probabile si ha $\alpha\log\pi_\alpha(s\mid k)\to0$, e la~\eqref{eq:score} si riduce a $D_0(k,s)$.
Con campi congelati il costo $\lnq b_k$ si paga comunque, prima o dopo, e il punteggio dipende solo dall'effetto della mossa sulle altre variabili.
Per questo motivo teniamo il bias di BSP come filtro dei candidati.
\end{osservazione}

La somma in~\eqref{eq:D} corre su tutte le variabili libere, ma riceve contributi soprattutto dalle variabili vicine a $k$.
Il Corollario~\ref{cor:tree} e la Proposizione~\ref{prop:remainder} suggeriscono questa località, ma essa è un'ipotesi da misurare.
Se la località vale, $D_\alpha$ è di ordine uno e non si diluisce come il secondo passo letterale.

\subsection{L'algoritmo SQ$(M,\alpha,q)$}

La variante cambia solo la scelta delle variabili da fissare in un passo di decimazione.
Sia $B=\max(1,\lfloor fN_t\rfloor)$ la dimensione del batch di BSP.
\begin{enumerate}
\item Si prendono come candidati le prime $M$ variabili libere nell'ordine di BSP. Se $M\le B$ l'algoritmo è BSP.
\item Per ogni candidato $k$ con il valore più probabile $s$ si esegue una sonda: si fissa $x_k=s$, si applica la propagazione unitaria e si riconverge SP su una copia della formula. Se la sonda trova una contraddizione o non converge, il punteggio vale $-\infty$. Altrimenti si calcola $V_\alpha(F^{(k,s)})$.
\item Il punteggio è il membro destro della~\eqref{eq:score}.
\item Per $\alpha=0$ si fissano i $B$ candidati con il punteggio più alto. Per $\alpha>0$ si estraggono $B$ candidati senza reinserimento con probabilità proporzionali a $e^{\text{punteggio}/\alpha}$, con il metodo Gumbel-top-$k$~\citep{kool2019stochastic}.
\item I passi di backtracking restano quelli di BSP.
\end{enumerate}
Per $M\le B$ si ritrova BSP.
Nel codice abbiamo verificato che questo limite, e il caso in cui le sonde vengono eseguite ma l'ordine di BSP viene mantenuto, riproducono BSP passo per passo.
Il costo è di $M$ riconvergenze di SP per ogni passo di decimazione.
Il look-ahead a un passo dell'Osservazione~\ref{oss:one-step} è il caso in cui il punteggio è la complessità dopo la sonda.

SAC stima $Q_\alpha$ con una rete e con l'apprendimento.
Qui il modello della dinamica è noto, perché è SP stessa, e $Q_\alpha$ si calcola con una sonda.
La famiglia di algoritmi più vicina è quindi la pianificazione a massima entropia~\citep{xiao2019maximum}, più che l'actor-critic.

\section{Diagrammi dei limiti}
\label{sec:diagrams}

Questa sezione riassume le derivazioni in quattro diagrammi.
Ogni freccia è un limite o una scelta di parametri, e l'etichetta dice quale.
Una freccia tratteggiata indica un limite formale, senza un significato fisico diretto.
Dopo ogni diagramma diciamo quali cammini commutano, cioè portano allo stesso oggetto.

\subsection{Potenziali 1RSB}

\begin{equation*}
\begin{tikzcd}[column sep=4.2em,row sep=3.2em]
& \Xi(\beta,m) \arrow[dl,"{\beta\to\infty,\ m\ \text{fisso}}"'] \arrow[dr,"{\beta\to\infty,\ \beta m=y}"] & \\
\varphi(m)\ \text{(entropico)} \arrow[d,"m=1"'] \arrow[dr,"m\to0"'] & & G(y)=-y\Phi(y)\ \text{(SP-}y) \arrow[dl,"y\to\infty"] \arrow[d,dashed,"y\to0"] \\
\text{BP sulle soluzioni} & \Sig_{\mathrm{SP}}\ \text{(SP)} & \sup_e\Sig_e(e)
\end{tikzcd}
\end{equation*}

Il cammino di sinistra conta i cluster di soluzioni pesati con $|C|^m$ e poi prende $m\to0$.
Il cammino di destra pesa gli stati con $e^{-yE_\alpha}$ e poi prende $y\to\infty$.
I due cammini commutano nella fase SAT, nel senso della Sezione~\ref{sec:notation}: le equazioni entropiche a $m\to0$, ristrette ai messaggi congelati, sono le equazioni SP.
Il vertice $m=1$ è BP sulla misura uniforme delle soluzioni.
La freccia $y\to0$ è tratteggiata perché il ramo a $y$ piccolo esce dalla regione di stabilità della 1RSB.
Sul lato sinistro l'ordine di Rényi dei pesi dei cluster è $m$ (Proposizione~\ref{prop:renyi-m}), e per $\alpha_c<\alpha<\alpha_s$ tutti gli ordini $m\ge m^*$ danno $h_m=0$.
Per $\alpha\to\alpha_s$ si ha $m^*\to0$, e la parte utile del lato sinistro si riduce al vertice SP.

\subsection{Semianelli termodinamici}

\begin{equation*}
\begin{tikzcd}[column sep=5.2em,row sep=3.6em]
\oplus_{\KL(\cdot\|\pi),T} \arrow[r,"\pi\ \text{uniforme}"] & \oplus_{\Sh,T} \arrow[d,"T\to0"] & \oplus_{S_q,T,q} \arrow[l,"q\to1"'] \arrow[dl,"{T\to0,\ q\le1}"'] \arrow[d,"{T\to0,\ q>1}"] \\
\oplus_{\Ry_\alpha,T} \arrow[ur,"\alpha\to1"] \arrow[r,"T\to0"'] & \min\ \text{(tropicale)} & M_q\ \text{(media di potenze)} \arrow[l,"q\to1^+"]
\end{tikzcd}
\end{equation*}

Qui $\oplus_{S_q,T,q}$ è la somma di Tsallis deformata~\eqref{eq:tsallis-op}, e $M_q$ è la media di potenze~\eqref{eq:power-mean}.
La freccia da $\oplus_{\KL(\cdot\|\pi),T}$ a $\oplus_{\Sh,T}$ vale a meno della costante $T\log n$.
Tutti i quadrati commutano come limiti iterati.
L'unico vertice non tropicale è $M_q$: a $q>1$ fisso il limite $T\to0$ della somma di Tsallis non è il minimo (Corollario~\ref{cor:tsallis-T0}).
Fra le operazioni del diagramma, sono associative e commutative solo $\oplus_{\Sh,T}$, il minimo e, con il prodotto deformato di MT, $\oplus_{S_q,T,q}$.
L'operazione $\oplus_{\Ry_\alpha,T}$ è commutativa ma non associativa.
L'operazione $\oplus_{\KL(\cdot\|\pi),T}$ con $\pi$ non uniforme non è commutativa.

\subsection{Dai semianelli ai livelli di SP}

\begin{equation*}
\begin{tikzcd}[column sep=3.4em,row sep=3.4em]
\mathcal Z_q(y),\ q-1=\kappa/N \arrow[r,"\text{Cor.~\ref{cor:yeff}}"] \arrow[d,"\kappa\to0"'] & \text{SP-}y\ \text{a}\ y_{\mathrm{eff}} \arrow[d,"y_{\mathrm{eff}}\to y"] & \\
\oplus_{\Sh,1/y}\{E_\alpha\}=N\Phi(y) \arrow[r,"\eqref{eq:reweight}"] & \text{SP-}y:\ \oplus_{\KL(\cdot\|P_0),1/y}\{\Delta E(\omega)\} \arrow[r,"y\to\infty"] & \text{SP} \\
\mathcal Z_q(y),\ q>1\ \text{fisso} \arrow[u,dashed,"{=\lim_{y\to0}}"'] & & \mathcal Z_q(y),\ q<1\ \text{fisso} \arrow[u,"{=\lim_{y\to\infty}}"']
\end{tikzcd}
\end{equation*}

La riga centrale è la derivazione della Sezione~\ref{sec:spy}.
La somma di Shannon sulle energie degli stati genera il reweighting di cavità, e il reweighting dà le somme di settore con entropia relativa.
La riga alta è la deformazione di Tsallis in scala $1/N$: essa dà di nuovo SP-$y$, con un $y$ efficace.
La riga bassa contiene le deformazioni di Tsallis a $q$ fisso, che degenerano nei due estremi dell'asse $y$ (Proposizione~\ref{prop:tsallis-regimes}).
Le frecce della riga bassa sono uguaglianze fra potenziali: con $q>1$ fisso si ottiene il potenziale del limite $y\to0$, con $q<1$ fisso quello del limite $y\to\infty$.

\subsection{Algoritmi}

\begin{equation*}
\begin{tikzcd}[column sep=3.6em,row sep=3.4em]
\text{SP-Y}(y,r),\ b_{\mathrm{back}} \arrow[r,"y\to\infty"] \arrow[d,"r=0"'] & \text{BSP}(r),\ b_{\mathrm{back}}|_{y=\infty} \arrow[d,"r=0"] & \text{BSP}(r),\ b_k \arrow[dl,"r=0"] \\
\text{SP-}y\ \text{e decimazione} \arrow[r,"y\to\infty"'] & \text{SID} &
\end{tikzcd}
\end{equation*}

Il quadrato a sinistra commuta.
Il vertice in alto al centro non è BSP: è BSP con l'indice~\eqref{eq:bkz-index} a $y=\infty$ al posto del bias~\eqref{eq:bias}.
BSP e questa variante differiscono solo nella regola di rilascio, e coincidono per $r=0$.
La Sezione~\ref{sec:response} dà una terza regola di rilascio, $I_k$ in~\eqref{eq:Ik}, che si colloca nello stesso vertice.

\begin{equation*}
\begin{tikzcd}[column sep=4.4em,row sep=3.4em]
& \text{BP sulle soluzioni} & \\
\text{BP}(\omega_o,\omega_*) \arrow[r,"\omega_o+\omega_*=1"] \arrow[d,"{\omega_o=0,\ \omega_*=e^\gamma}"'] & \text{SP}(\rho),\ \rho=\omega_* \arrow[r,"\rho=1"] \arrow[u,"\rho=0"'] & \text{SP} \\
\text{tilt}\ \gamma\ \text{(messaggi completi)} \arrow[urr,"\gamma=0"'] & &
\end{tikzcd}
\end{equation*}

Questo diagramma è il piano di MMW.
La retta $\omega_o+\omega_*=1$ è la famiglia SP$(\rho)$, e su di essa i messaggi si chiudono su una sola survey.
La retta $\omega_o=0$ è il tilt $\gamma$ verso gli stati $*$.
Essa incontra la retta di SP$(\rho)$ solo nel punto SP, per la Proposizione~\ref{prop:gamma}.

\section{Discussione e problemi aperti}
\label{sec:discussion}

\subsection{Cosa è noto e cosa aggiungiamo}

La Tabella~\ref{tab:known} separa i risultati noti dalle osservazioni di questo lavoro.

\begin{table}[t]
\centering
\small
\begin{tabular}{@{}p{0.55\linewidth}p{0.38\linewidth}@{}}
\toprule
Enunciato & Fonte \\
\midrule
Potenziale 1RSB $\Xi(\beta,m)$, limiti entropico ed energetico & \citep{monasson1995structural,mezard2003cavity,mezard2002random} \\
SP, SP-$y$ e funzionale di complessità & \citep{mezard2002analytic,mezard2002random,braunstein2005survey} \\
SP come BP su uno spazio esteso, famiglia SP$(\rho)$ & \citep{braunstein2004surveybp,maneva2007survey} \\
Bias $b_i$ come fattore sul numero di cluster; BSP & \citep{marino2016bsp} \\
SP-$y$ con backtracking e indice con segno & \citep{battaglia2004minimizing} \\
Semianelli termodinamici, Shannon, Rényi, Tsallis, KL & \citep{marcolli2014thermodynamic} \\
Ensemble canonici generalizzati & \citep{costeniuc2005generalized,costeniuc2006generalized} \\
\midrule
Stazionarietà del funzionale completo, non della forma compatta (Lemma~\ref{lem:stationary}) & dimostrazione diretta nostra di un fatto generale~\citep{yedidia2005constructing} \\
Due livelli della somma di Shannon; $y$ locale uguale a $y$ dei cluster & nostro \\
Settori di SP-$y$ come somme con entropia relativa & nostro \\
Ordine di Rényi uguale a $m$; spettro multifrattale; $h_m=0$ per $m\ge m^*$ & nostro \\
Forma chiusa e limite $T\to0$ della somma di Tsallis deformata & nostro, con~\citep{curado1991generalized} \\
Tsallis sui cluster: tre regimi, $y_{\mathrm{eff}}$ & nostro \\
Tilt $\gamma$: rottura dell'invariante di MMW & nostro, da~\citep{maneva2007survey} \\
Identità locale con segno e resto di secondo ordine & nostro; il caso del valore più probabile è in~\citep{marino2016bsp} \\
Equazioni di Bellman soft e politica di Gibbs & \citep{haarnoja2018soft} \\
Look-ahead soft a due passi, valore a campi congelati, ricompensa di Tsallis (Sezione~\ref{sec:softq}) & nostro; verificato in Lean~4 \\
\bottomrule
\end{tabular}
\caption{Risultati noti (sopra) e osservazioni di questo lavoro (sotto).}
\label{tab:known}
\end{table}

Il quadro dei semianelli non produce da solo un algoritmo nuovo.
A $T$ fisso il semianello di Shannon è isomorfo al semianello di BP, quindi ogni equazione di cavità si può scrivere in esso.
Il valore del quadro sta nei limiti e nelle deformazioni.
Esso dice quali deformazioni restano associative, quali degenerano, e a quale livello della teoria agisce ogni parametro.

\subsection{Problemi aperti}

\paragraph{Non concavità di $\Sig_e$.}
La deformazione di Tsallis in scala $1/N$ con $q<1$ dà stati nuovi solo se $\Sig_e(e)$ ha tratti non concavi.
Per saperlo si può calcolare $\Sig_e(e)$ con la dinamica di popolazione vicino ad $\alpha_s$, per $K=3$ e $K=4$.

\paragraph{Il ciclo su $y_{\mathrm{eff}}$.}
Il Corollario~\ref{cor:yeff} suggerisce un ciclo esterno che aggiorna $y_{\mathrm{eff}}$ con l'energia del funzionale di Bethe.
Resta da capire se il ciclo converge e se, con backtracking, cambia la soglia algoritmica.

\paragraph{Il tilt $\gamma$.}
BP sul modello di MMW con $\omega_o=0$ e $\omega_*=e^\gamma$ è ben definita, ma non si chiude su una survey.
Resta da confrontare questa misura con la chiusura approssimata del nostro codice.
Resta anche da capire se la misura aiuta la ricerca di cluster non congelati, come suggerito in~\citep{marino2016bsp}.

\paragraph{Il resto $R_{k,s}$ su grafi con cicli.}
La Proposizione~\ref{prop:remainder} dà l'ordine del resto ma non la sua grandezza.
Una stima tramite la suscettività di SP e la lunghezza dei cicli renderebbe quantitativo il confronto fra $I_k$, $b_k$ e $b_{\mathrm{back}}$.
Ricci-Tersenghi e Semerjian~\citep{ricci2009cavity} studiano la decimazione con il metodo della cavità, e il loro approccio è un punto di partenza.

\paragraph{Esperimenti.}
Le prove numeriche fatte finora sono troppo piccole.
Un confronto utile deve fissare il budget di calcolo, usare molte istanze per ogni $(K,\alpha,N)$, e riportare la frazione di istanze risolte con il suo errore.
Le regole di rilascio da confrontare sono $b_k$, $b_{\mathrm{back}}$ e $I_k$, a parità di $r$ e $f$.

\section{Conclusioni}

Abbiamo scritto le deformazioni note di SP nella lingua dei semianelli termodinamici e abbiamo derivato ogni passo.
La somma di Shannon agisce su due livelli, e il livello dei cluster fissa il reweighting locale.
L'ordine di Rényi dei pesi dei cluster è il parametro di Parisi, e vicino alla soglia la parte utile dello spettro multifrattale si riduce a SP.
La deformazione di Tsallis con $q$ fisso degenera, mentre la deformazione in scala $1/N$ è un ensemble canonico generalizzato che si riduce a SP-$y$ con un $y$ efficace.
L'identità locale della complessità giustifica il bias di BSP per il fissaggio e dà una regola con segno per il rilascio, con un resto di secondo ordine.
I diagrammi dei limiti raccolgono questi fatti e mostrano quali varianti sono nuove e quali sono riparametrizzazioni di algoritmi noti.


\appendix
\input{sections/appendix_bsp}

\bibliography{thermosp}
\bibliographystyle{iclr2027_conference}

\end{document}
